%
%Academic CV LaTeX Template
% Author: Dubasi Pavan Kumar
% LinkedIn: https://www.linkedin.com/in/im-pavankumar/
% License: MIT
%
% For errors, suggestions, or improvements, please contact:
% Email: pavankumard.pg19.ma@nitp.ac.in
%
%
% -----------------------------------------------------------------------------
% THIS FILE IS STYLING AND MACROS ONLY, SHARED BY EVERY DOCUMENT UNDER cv/.
%
% Every fact printed below comes from `content/cv.yaml` via the generated macros
% in `cv/generated/cv-data.tex` (produced by `node scripts/build-cv-data.mjs`).
% Publications and talks come straight from `content/publications.bib` and
% `content/talks.bib` through the biblatex refsections at the end of this file.
%
% Packages, page geometry, fonts, colours, the biblatex setup and the section
% macros live here. Every document under cv/ opens with %
%Academic CV LaTeX Template
% Author: Dubasi Pavan Kumar
% LinkedIn: https://www.linkedin.com/in/im-pavankumar/
% License: MIT
%
% For errors, suggestions, or improvements, please contact:
% Email: pavankumard.pg19.ma@nitp.ac.in
%
%
% -----------------------------------------------------------------------------
% THIS FILE IS STYLING AND MACROS ONLY, SHARED BY EVERY DOCUMENT UNDER cv/.
%
% Every fact printed below comes from `content/cv.yaml` via the generated macros
% in `cv/generated/cv-data.tex` (produced by `node scripts/build-cv-data.mjs`).
% Publications and talks come straight from `content/publications.bib` and
% `content/talks.bib` through the biblatex refsections at the end of this file.
%
% Packages, page geometry, fonts, colours, the biblatex setup and the section
% macros live here. Every document under cv/ opens with %
%Academic CV LaTeX Template
% Author: Dubasi Pavan Kumar
% LinkedIn: https://www.linkedin.com/in/im-pavankumar/
% License: MIT
%
% For errors, suggestions, or improvements, please contact:
% Email: pavankumard.pg19.ma@nitp.ac.in
%
%
% -----------------------------------------------------------------------------
% THIS FILE IS STYLING AND MACROS ONLY, SHARED BY EVERY DOCUMENT UNDER cv/.
%
% Every fact printed below comes from `content/cv.yaml` via the generated macros
% in `cv/generated/cv-data.tex` (produced by `node scripts/build-cv-data.mjs`).
% Publications and talks come straight from `content/publications.bib` and
% `content/talks.bib` through the biblatex refsections at the end of this file.
%
% Packages, page geometry, fonts, colours, the biblatex setup and the section
% macros live here. Every document under cv/ opens with %
%Academic CV LaTeX Template
% Author: Dubasi Pavan Kumar
% LinkedIn: https://www.linkedin.com/in/im-pavankumar/
% License: MIT
%
% For errors, suggestions, or improvements, please contact:
% Email: pavankumard.pg19.ma@nitp.ac.in
%
%
% -----------------------------------------------------------------------------
% THIS FILE IS STYLING AND MACROS ONLY, SHARED BY EVERY DOCUMENT UNDER cv/.
%
% Every fact printed below comes from `content/cv.yaml` via the generated macros
% in `cv/generated/cv-data.tex` (produced by `node scripts/build-cv-data.mjs`).
% Publications and talks come straight from `content/publications.bib` and
% `content/talks.bib` through the biblatex refsections at the end of this file.
%
% Packages, page geometry, fonts, colours, the biblatex setup and the section
% macros live here. Every document under cv/ opens with \input{preamble.tex},
% so a change to the printed design is made once and reaches all of them.
%
% To change what a CV says, edit content/cv.yaml and regenerate. To change how
% it looks, edit this file. To change WHICH of it one document prints, edit that
% document's body - cv/cv.tex, cv/short.tex or cv/teaching.tex.
%
% WHICH sections print is a record fact, not a layout one. The `\cvpart` lines in
% a document body places the sections whose heading or setting that document
% curates; `\cvAutoSections` then prints every remaining section of
% content/cv.yaml, in the order the record writes them, under the heading its key
% spells out. Adding a section to the record therefore needs no edit here. Give
% a section `printed: false` in content/cv.yaml to keep it off the PDF - that
% holds for the curated ones below too - and `heading:` to name an uncurated one
% something other than its key.
%
% Build:  latexmk -xelatex -cd cv/cv.tex      (xelatex is required: academicons)
%         the same for cv/short.tex and cv/teaching.tex
% -----------------------------------------------------------------------------



\documentclass[a4paper,11pt]{article}

% Package imports
\usepackage{latexsym}
\usepackage{xcolor}
\usepackage{float}
\usepackage{ragged2e}
\usepackage[empty]{fullpage}
\usepackage{wrapfig}
\usepackage{lipsum}
\usepackage{booktabs}   % high-quality horizontal rules (no vertical rules)
\usepackage{array}      % better column types
\usepackage{xparse}
\usepackage{tabularx}
\usepackage{titlesec}
\usepackage{geometry}
\usepackage{marvosym}
\usepackage{verbatim}
\usepackage{enumitem}
\usepackage{fancyhdr}
\usepackage{multicol}
\usepackage{graphicx}
\usepackage{fontspec}
\setmainfont{TeX Gyre Pagella}
\usepackage{fontawesome5}
\usepackage{academicons} % <-- NEW (Scholar, ORCID, etc.)
\usepackage{eurosym} % <-- NEW (€)

% Column helpers (tight but readable)
\newcolumntype{P}[1]{>{\raggedright\arraybackslash}p{#1}}
\newcolumntype{Y}{>{\raggedright\arraybackslash}X}

% A reusable CV table environment (no floats, correct spacing)
\NewDocumentEnvironment{cvtable}{m +b}{%
  \vspace{0.25em}
  \small
  \setlength{\tabcolsep}{5pt}
  \renewcommand{\arraystretch}{1.15}
  \begin{tabularx}{\linewidth}{@{}#1@{}}
  #2
  \end{tabularx}
  \vspace{0.25em}
}{}

% Color definitions
\definecolor{darkblue}{RGB}{0,0,139}

% Page layout
\setlength{\multicolsep}{0pt}
\pagestyle{fancy}
\fancyhf{} % clear all header and footer fields
\fancyfoot{}
\renewcommand{\headrulewidth}{0pt}
\renewcommand{\footrulewidth}{0pt}
\geometry{left=1.4cm, top=0.8cm, right=1.2cm, bottom=1cm}
\setlength{\footskip}{5pt} % Addressing fancyhdr warning

% Hyperlink setup (moved after fancyhdr to address warning)
\usepackage[hidelinks]{hyperref}
\hypersetup{
    colorlinks=true,
    linkcolor=darkblue,
    filecolor=darkblue,
    urlcolor=darkblue,
}

% Custom box settings
\usepackage[most]{tcolorbox}
\tcbset{
    frame code={},
    center title,
    left=0pt,
    right=0pt,
    top=0pt,
    bottom=0pt,
    colback=gray!20,
    colframe=white,
    width=\dimexpr\textwidth\relax,
    enlarge left by=-2mm,
    boxsep=4pt,
    arc=0pt,outer arc=0pt,
}

% URL style
\urlstyle{same}

% Text alignment
\raggedright
\setlength{\tabcolsep}{0in}

% Section formatting
\titleformat{\section}{
  \vspace{-4pt}\scshape\raggedright\large
}{}{0em}{}[\color{black}\titlerule \vspace{-7pt}]

% Custom commands
\newcommand{\resumeItem}[2]{
  \item{
    \textbf{#1}{\hspace{0.5mm}#2 \vspace{-0.5mm}}
  }
}

\newcommand{\resumePOR}[3]{
\vspace{0.5mm}\item
    \begin{tabular*}{0.97\textwidth}[t]{l@{\extracolsep{\fill}}r}
        \textbf{#1}\hspace{0.3mm}#2 & \textit{\small{#3}}
    \end{tabular*}
    \vspace{-2mm}
}

% One entry, every section: what it was, where it was, when, and the place.
% `content/cv.yaml` has exactly one entry shape and this is how it is set.
\newcommand{\cventry}[4]{
\vspace{0.5mm}\item
    \begin{tabularx}{0.98\textwidth}[t]{@{}Y >{\raggedleft\arraybackslash}p{4.8cm}@{}}
        \textbf{#1} & \textit{\footnotesize{#3}} \\
        \textit{\footnotesize{#2}} &  \footnotesize{#4}\\
    \end{tabularx}
    \vspace{-2.4mm}
}

% An empty generated macro, to test the optional prose fields against.
% NOT `\ifx\cvFocus\empty`: \newcommand makes its macros \long and \empty is
% not, so \ifx compares false for two macros that both hold nothing, and the
% heading prints over an empty field. This one is made the same way they are.
\newcommand{\cvempty}{}

% Separates the paragraphs of a section's `note:`. Spacing is layout, so the
% amount of it lives here and the generator only says where the break goes.
\newcommand{\cvnotesep}{\par\vspace{1.2mm}}

% Record that this file lays out one section key itself, so \cvautopart below
% does not print it a second time from the record's own section sequence.
% \gdef, not \def: \cvpartflush sets its entries inside a group.
\newcommand{\cvlaidout}[1]{\expandafter\gdef\csname cvlaidout@#1\endcsname{}}

% Make the five generated macros this layout may mention safe even when
% content/cv.yaml has no such section, and default the section to printing.
% (`Header` is the sixth generated macro, but this layout supplies its own table
% headings.) `\cvpart` does this for you; a hand-written section block must call it
% itself.
\newcommand{\cvdeclare}[1]{%
  \cvlaidout{#1}%
  % Printing is the default: only `printed: false` in content/cv.yaml generates a
  % \cv<Key>Printed of 0, and this \providecommand leaves that definition alone.
  \expandafter\providecommand\csname cv#1Printed\endcsname{1}%
  \expandafter\providecommand\csname cv#1Count\endcsname{0}%
  \expandafter\providecommand\csname cv#1\endcsname{}%
  \expandafter\providecommand\csname cv#1Note\endcsname{}%
  \expandafter\providecommand\csname cv#1Rows\endcsname{}%
  \expandafter\providecommand\csname cv#1Inline\endcsname{}%
}

% The same, for a bibliography section: `publications:`/`talks:` in cv.yaml
% generate \cv<Name>Key, \cv<Name>Sections and \cv<Name>Count (the number of
% entries in the matching .bib) instead of the five above. An adopter who has
% declared neither - which is everyone on day one - would otherwise meet a raw
% "Undefined control sequence" that stops xelatex at its interactive prompt.
% Empty is the documented behaviour: the heading prints over an empty
% bibliography and README.md tells them to delete the block.
\newcommand{\cvdeclarebib}[1]{%
  \expandafter\providecommand\csname cv#1Count\endcsname{0}%
  \expandafter\providecommand\csname cv#1Key\endcsname{}%
  \expandafter\providecommand\csname cv#1Sections\endcsname{}%
}

% How many entries of a section a document prints. The generated \cv<Name>
% guards each of its entries with \ifnum<n>>\cvmax, so this is the cutoff.
%
% This is the layout's call, and it is a different question from the record's
% `printed:`. WHICH sections print is a record fact; HOW MANY entries of one a
% document prints is a page budget belonging to that document - the same section
% is worth three entries in a short CV and all of them in the full one - so the
% number lives here and content/cv.yaml never states it. A count high enough
% that no record reaches it is the default, which keeps the full CV whole.
%
% Curation is neither: "my five best papers" is a judgement about the work, so
% it is marked on the fact itself - see the `selected` keyword in
% content/publications.bib and how cv/short.tex filters on it.
\newcommand{\cvmax}{99999}

% Print one section of content/cv.yaml. #2 is the heading, #3 the section's
% macro name (`Awards` for the `awards:` key), and the optional #1 is how many
% of its entries to print. A section the file does not have, or gives
% `printed: false`, prints nothing at all - no heading, no gap - so this line
% can stay. Reorder the CV by reordering these lines in the document body.
%
% The whole thing is one group so that \cvmax is restored afterwards: a
% \cvpart[3] must not silently truncate the section written after it.
\newcommand{\cvpart}[3][99999]{{%
  \renewcommand{\cvmax}{#1}%
  \cvdeclare{#3}%
  \ifnum\csname cv#3Count\endcsname>0
  \ifnum\csname cv#3Printed\endcsname>0
  \section{\textbf{#2}}
  \vspace{-0.4mm}
  \resumeSubHeadingListStart
  \csname cv#3\endcsname
  \resumeSubHeadingListEnd
  \vspace{-6mm}
  \fi
  \fi
}}

% One section of the record's own section sequence: `\cvpart`, unless this file
% has already laid that key out by hand above. Every line of \cvAutoSections is
% one of these, so an adopter's new section prints with no edit here, and the
% ones curated below keep their curated heading, order and setting.
\newcommand{\cvautopart}[2]{%
  \expandafter\ifx\csname cvlaidout@#2\endcsname\relax\cvpart{#1}{#2}\fi%
}

% `\cvpart`, with the second line of every entry set the other way. Which
% sections use which is a layout choice and lives here; cv.yaml says nothing
% about it.
\newcommand{\cvpartflush}[3][99999]{{%
  \renewcommand{\cventry}[4]{%
    \vspace{0.5mm}\item
    \begin{tabularx}{0.98\textwidth}[t]{@{}Y >{\raggedleft\arraybackslash}p{4.8cm}@{}}
        \textbf{##1} & \textit{\footnotesize{##3}} \\
        \footnotesize{\textit{##2}} & \footnotesize{##4}
    \end{tabularx}
    \vspace{-2.4mm}
  }%
  \cvpart[#1]{#2}{#3}}}

% A table hanging under an entry: `rows:` in cv.yaml. #1 is a generated
% equal-width column specification, #2 the generated headings and #3 the rows.
% Any number of consistently shaped columns therefore works without editing
% this layout. The generator refuses rows whose key order disagrees.
\newcommand{\cvcourses}[3]{%
\par
\begin{cvtable}{#1}
\toprule
#2
\midrule
#3
\bottomrule
\end{cvtable}
}

\newcommand{\resumeSubItem}[2]{\resumeItem{#1}{#2}\vspace{-4pt}}

\renewcommand{\labelitemi}{$\vcenter{\hbox{\tiny$\bullet$}}$}
\renewcommand{\labelitemii}{$\vcenter{\hbox{\tiny$\circ$}}$}

\newcommand{\resumeSubHeadingListStart}{\begin{itemize}[leftmargin=*,labelsep=1mm]}
\newcommand{\resumeHeadingSkillStart}{\begin{itemize}[leftmargin=*,itemsep=1.7mm, rightmargin=2ex]}
\newcommand{\resumeItemListStart}{\begin{itemize}[leftmargin=*,labelsep=1mm,itemsep=0.5mm]}

\newcommand{\resumeSubHeadingListEnd}{\end{itemize}\vspace{2mm}}
\newcommand{\resumeHeadingSkillEnd}{\end{itemize}\vspace{-2mm}}
\newcommand{\resumeItemListEnd}{\end{itemize}\vspace{-2mm}}
\newcommand{\cvsection}[1]{%
\vspace{2mm}
\begin{tcolorbox}
    \textbf{\large #1}
\end{tcolorbox}
    \vspace{-4mm}
}

\newcolumntype{L}{>{\raggedright\arraybackslash}X}%
\newcolumntype{R}{>{\raggedleft\arraybackslash}X}%
\newcolumntype{C}{>{\centering\arraybackslash}X}%

% Commands for icon sizing and positioning
\newcommand{\socialicon}[1]{\raisebox{-0.05em}{\resizebox{!}{1em}{#1}}}
\newcommand{\ieeeicon}[1]{\raisebox{-0.3em}{\resizebox{!}{1.3em}{#1}}}

% Profile icon registry. Known compact account IDs select their service icon;
% arbitrary `{label, url}` links use \cviconlink.
\newcommand{\cviconscholar}{\socialicon{\aiGoogleScholar}}
\newcommand{\cviconorcid}{\socialicon{\aiOrcid}}
\newcommand{\cviconlinkedin}{\socialicon{\faLinkedin}}
\newcommand{\cvicongithub}{\socialicon{\faGithub}}
\newcommand{\cviconlink}{\socialicon{\faLink}}

% Font options
\newcommand{\headerfonti}{\fontfamily{phv}\selectfont} % Helvetica-like (similar to Arial/Calibri)
\newcommand{\headerfontii}{\fontfamily{ptm}\selectfont} % Times-like (similar to Times New Roman)
\newcommand{\headerfontiii}{\rmfamily} % TeX Gyre Pagella, set above with Unicode support
\newcommand{\headerfontiv}{\fontfamily{pbk}\selectfont} % Bookman (readable serif)
\newcommand{\headerfontv}{\fontfamily{pag}\selectfont} % Avant Garde-like (similar to Trebuchet MS)
\newcommand{\headerfontvi}{\fontfamily{cmss}\selectfont} % Computer Modern Sans Serif
\newcommand{\headerfontvii}{\fontfamily{qhv}\selectfont} % Quasi-Helvetica (another Arial/Calibri alternative)
\newcommand{\headerfontviii}{\fontfamily{qpl}\selectfont} % Quasi-Palatino (another elegant serif option)
\newcommand{\headerfontix}{\fontfamily{qtm}\selectfont} % Quasi-Times (another Times New Roman alternative)
\newcommand{\headerfontx}{\fontfamily{bch}\selectfont} % Charter (clean serif font)

% --------------------------------
% BIBLATEX (publications + talks)
% --------------------------------
\usepackage[
  backend=biber,
  style=numeric,
  sorting=ydnt,
  defernumbers=true,
  maxbibnames=99,
  giveninits=true,
  doi=true,
  url=true,
  isbn=false
]{biblatex}

% Two bib databases (kept separate via refsection below).
% Publications are the repository-wide canonical list, shared with the website.
\addbibresource{../content/publications.bib}
\addbibresource{../content/talks.bib}

% The shared bibliography may carry rendering fields used only by the website. They are not part of the CV, and abstracts can contain raw % and &
% that would break the LaTeX pass, so drop them during biber processing instead
% of altering the .bib the website reads.
\DeclareSourcemap{
  \maps[datatype=bibtex]{
    \map{
      \step[fieldset=abstract, null]
      \step[fieldset=selected, null]
      \step[fieldset=preview, null]
      \step[fieldset=abbr, null]
      \step[fieldset=bibtex_show, null]
      \step[fieldset=html, null]
      \step[fieldset=pdf, null]
      \step[fieldset=code, null]
      \step[fieldset=altmetric, null]
      \step[fieldset=dimensions, null]
    }
    % DBLP's @misc artifact records carry an UNFILLED placeholder,
    % `note = {Accessed on YYYY-MM-DD.}`, which biblatex would print verbatim as
    % the venue of a released software package. The website already skips it
    % (`VENUE_FIELDS` in web/src/lib/record.ts prefers `publisher`); this is the
    % same fix on the printed side, so declaring a `types: [misc]` section in
    % content/cv.yaml renders correctly instead of printing the placeholder.
    \map{
      \step[fieldsource=note, match=\regexp{\AAccessed\son\sYYYY-MM-DD\.?\Z}, final]
      \step[fieldset=note, null]
    }
  }
}

% Compact bib
\setlength\bibitemsep{0.25\baselineskip}
\renewcommand*{\bibfont}{\footnotesize}

% Bib headings matching your tidy style
\defbibheading{bibsubheading}[\refname]{%
  \vspace{0.8mm}\textbf{#1}\par\vspace{-0.3mm}
}

% No filters here. Which entry type belongs under which heading is a curated
% opinion, so it is declared in `content/cv.yaml` under `publications:` and
% `talks:` and read by the website too. build-cv-data.mjs turns each declared
% section into its \defbibfilter and its \printbibliography, and gives this file
% \cvPublicationsKey / \cvPublicationsSections and \cvTalksKey /
% \cvTalksSections, plus \cvPublicationsCount / \cvTalksCount - how many entries
% the .bib holds, which guards the \refsection below.
% Adding, reordering or renaming a section is a cv.yaml edit.

% --------------------------------
% CONTENT (generated from content/cv.yaml)
% --------------------------------
\input{generated/cv-data.tex}

,
% so a change to the printed design is made once and reaches all of them.
%
% To change what a CV says, edit content/cv.yaml and regenerate. To change how
% it looks, edit this file. To change WHICH of it one document prints, edit that
% document's body - cv/cv.tex, cv/short.tex or cv/teaching.tex.
%
% WHICH sections print is a record fact, not a layout one. The `\cvpart` lines in
% a document body places the sections whose heading or setting that document
% curates; `\cvAutoSections` then prints every remaining section of
% content/cv.yaml, in the order the record writes them, under the heading its key
% spells out. Adding a section to the record therefore needs no edit here. Give
% a section `printed: false` in content/cv.yaml to keep it off the PDF - that
% holds for the curated ones below too - and `heading:` to name an uncurated one
% something other than its key.
%
% Build:  latexmk -xelatex -cd cv/cv.tex      (xelatex is required: academicons)
%         the same for cv/short.tex and cv/teaching.tex
% -----------------------------------------------------------------------------



\documentclass[a4paper,11pt]{article}

% Package imports
\usepackage{latexsym}
\usepackage{xcolor}
\usepackage{float}
\usepackage{ragged2e}
\usepackage[empty]{fullpage}
\usepackage{wrapfig}
\usepackage{lipsum}
\usepackage{booktabs}   % high-quality horizontal rules (no vertical rules)
\usepackage{array}      % better column types
\usepackage{xparse}
\usepackage{tabularx}
\usepackage{titlesec}
\usepackage{geometry}
\usepackage{marvosym}
\usepackage{verbatim}
\usepackage{enumitem}
\usepackage{fancyhdr}
\usepackage{multicol}
\usepackage{graphicx}
\usepackage{fontspec}
\setmainfont{TeX Gyre Pagella}
\usepackage{fontawesome5}
\usepackage{academicons} % <-- NEW (Scholar, ORCID, etc.)
\usepackage{eurosym} % <-- NEW (€)

% Column helpers (tight but readable)
\newcolumntype{P}[1]{>{\raggedright\arraybackslash}p{#1}}
\newcolumntype{Y}{>{\raggedright\arraybackslash}X}

% A reusable CV table environment (no floats, correct spacing)
\NewDocumentEnvironment{cvtable}{m +b}{%
  \vspace{0.25em}
  \small
  \setlength{\tabcolsep}{5pt}
  \renewcommand{\arraystretch}{1.15}
  \begin{tabularx}{\linewidth}{@{}#1@{}}
  #2
  \end{tabularx}
  \vspace{0.25em}
}{}

% Color definitions
\definecolor{darkblue}{RGB}{0,0,139}

% Page layout
\setlength{\multicolsep}{0pt}
\pagestyle{fancy}
\fancyhf{} % clear all header and footer fields
\fancyfoot{}
\renewcommand{\headrulewidth}{0pt}
\renewcommand{\footrulewidth}{0pt}
\geometry{left=1.4cm, top=0.8cm, right=1.2cm, bottom=1cm}
\setlength{\footskip}{5pt} % Addressing fancyhdr warning

% Hyperlink setup (moved after fancyhdr to address warning)
\usepackage[hidelinks]{hyperref}
\hypersetup{
    colorlinks=true,
    linkcolor=darkblue,
    filecolor=darkblue,
    urlcolor=darkblue,
}

% Custom box settings
\usepackage[most]{tcolorbox}
\tcbset{
    frame code={},
    center title,
    left=0pt,
    right=0pt,
    top=0pt,
    bottom=0pt,
    colback=gray!20,
    colframe=white,
    width=\dimexpr\textwidth\relax,
    enlarge left by=-2mm,
    boxsep=4pt,
    arc=0pt,outer arc=0pt,
}

% URL style
\urlstyle{same}

% Text alignment
\raggedright
\setlength{\tabcolsep}{0in}

% Section formatting
\titleformat{\section}{
  \vspace{-4pt}\scshape\raggedright\large
}{}{0em}{}[\color{black}\titlerule \vspace{-7pt}]

% Custom commands
\newcommand{\resumeItem}[2]{
  \item{
    \textbf{#1}{\hspace{0.5mm}#2 \vspace{-0.5mm}}
  }
}

\newcommand{\resumePOR}[3]{
\vspace{0.5mm}\item
    \begin{tabular*}{0.97\textwidth}[t]{l@{\extracolsep{\fill}}r}
        \textbf{#1}\hspace{0.3mm}#2 & \textit{\small{#3}}
    \end{tabular*}
    \vspace{-2mm}
}

% One entry, every section: what it was, where it was, when, and the place.
% `content/cv.yaml` has exactly one entry shape and this is how it is set.
\newcommand{\cventry}[4]{
\vspace{0.5mm}\item
    \begin{tabularx}{0.98\textwidth}[t]{@{}Y >{\raggedleft\arraybackslash}p{4.8cm}@{}}
        \textbf{#1} & \textit{\footnotesize{#3}} \\
        \textit{\footnotesize{#2}} &  \footnotesize{#4}\\
    \end{tabularx}
    \vspace{-2.4mm}
}

% An empty generated macro, to test the optional prose fields against.
% NOT `\ifx\cvFocus\empty`: \newcommand makes its macros \long and \empty is
% not, so \ifx compares false for two macros that both hold nothing, and the
% heading prints over an empty field. This one is made the same way they are.
\newcommand{\cvempty}{}

% Separates the paragraphs of a section's `note:`. Spacing is layout, so the
% amount of it lives here and the generator only says where the break goes.
\newcommand{\cvnotesep}{\par\vspace{1.2mm}}

% Record that this file lays out one section key itself, so \cvautopart below
% does not print it a second time from the record's own section sequence.
% \gdef, not \def: \cvpartflush sets its entries inside a group.
\newcommand{\cvlaidout}[1]{\expandafter\gdef\csname cvlaidout@#1\endcsname{}}

% Make the five generated macros this layout may mention safe even when
% content/cv.yaml has no such section, and default the section to printing.
% (`Header` is the sixth generated macro, but this layout supplies its own table
% headings.) `\cvpart` does this for you; a hand-written section block must call it
% itself.
\newcommand{\cvdeclare}[1]{%
  \cvlaidout{#1}%
  % Printing is the default: only `printed: false` in content/cv.yaml generates a
  % \cv<Key>Printed of 0, and this \providecommand leaves that definition alone.
  \expandafter\providecommand\csname cv#1Printed\endcsname{1}%
  \expandafter\providecommand\csname cv#1Count\endcsname{0}%
  \expandafter\providecommand\csname cv#1\endcsname{}%
  \expandafter\providecommand\csname cv#1Note\endcsname{}%
  \expandafter\providecommand\csname cv#1Rows\endcsname{}%
  \expandafter\providecommand\csname cv#1Inline\endcsname{}%
}

% The same, for a bibliography section: `publications:`/`talks:` in cv.yaml
% generate \cv<Name>Key, \cv<Name>Sections and \cv<Name>Count (the number of
% entries in the matching .bib) instead of the five above. An adopter who has
% declared neither - which is everyone on day one - would otherwise meet a raw
% "Undefined control sequence" that stops xelatex at its interactive prompt.
% Empty is the documented behaviour: the heading prints over an empty
% bibliography and README.md tells them to delete the block.
\newcommand{\cvdeclarebib}[1]{%
  \expandafter\providecommand\csname cv#1Count\endcsname{0}%
  \expandafter\providecommand\csname cv#1Key\endcsname{}%
  \expandafter\providecommand\csname cv#1Sections\endcsname{}%
}

% How many entries of a section a document prints. The generated \cv<Name>
% guards each of its entries with \ifnum<n>>\cvmax, so this is the cutoff.
%
% This is the layout's call, and it is a different question from the record's
% `printed:`. WHICH sections print is a record fact; HOW MANY entries of one a
% document prints is a page budget belonging to that document - the same section
% is worth three entries in a short CV and all of them in the full one - so the
% number lives here and content/cv.yaml never states it. A count high enough
% that no record reaches it is the default, which keeps the full CV whole.
%
% Curation is neither: "my five best papers" is a judgement about the work, so
% it is marked on the fact itself - see the `selected` keyword in
% content/publications.bib and how cv/short.tex filters on it.
\newcommand{\cvmax}{99999}

% Print one section of content/cv.yaml. #2 is the heading, #3 the section's
% macro name (`Awards` for the `awards:` key), and the optional #1 is how many
% of its entries to print. A section the file does not have, or gives
% `printed: false`, prints nothing at all - no heading, no gap - so this line
% can stay. Reorder the CV by reordering these lines in the document body.
%
% The whole thing is one group so that \cvmax is restored afterwards: a
% \cvpart[3] must not silently truncate the section written after it.
\newcommand{\cvpart}[3][99999]{{%
  \renewcommand{\cvmax}{#1}%
  \cvdeclare{#3}%
  \ifnum\csname cv#3Count\endcsname>0
  \ifnum\csname cv#3Printed\endcsname>0
  \section{\textbf{#2}}
  \vspace{-0.4mm}
  \resumeSubHeadingListStart
  \csname cv#3\endcsname
  \resumeSubHeadingListEnd
  \vspace{-6mm}
  \fi
  \fi
}}

% One section of the record's own section sequence: `\cvpart`, unless this file
% has already laid that key out by hand above. Every line of \cvAutoSections is
% one of these, so an adopter's new section prints with no edit here, and the
% ones curated below keep their curated heading, order and setting.
\newcommand{\cvautopart}[2]{%
  \expandafter\ifx\csname cvlaidout@#2\endcsname\relax\cvpart{#1}{#2}\fi%
}

% `\cvpart`, with the second line of every entry set the other way. Which
% sections use which is a layout choice and lives here; cv.yaml says nothing
% about it.
\newcommand{\cvpartflush}[3][99999]{{%
  \renewcommand{\cventry}[4]{%
    \vspace{0.5mm}\item
    \begin{tabularx}{0.98\textwidth}[t]{@{}Y >{\raggedleft\arraybackslash}p{4.8cm}@{}}
        \textbf{##1} & \textit{\footnotesize{##3}} \\
        \footnotesize{\textit{##2}} & \footnotesize{##4}
    \end{tabularx}
    \vspace{-2.4mm}
  }%
  \cvpart[#1]{#2}{#3}}}

% A table hanging under an entry: `rows:` in cv.yaml. #1 is a generated
% equal-width column specification, #2 the generated headings and #3 the rows.
% Any number of consistently shaped columns therefore works without editing
% this layout. The generator refuses rows whose key order disagrees.
\newcommand{\cvcourses}[3]{%
\par
\begin{cvtable}{#1}
\toprule
#2
\midrule
#3
\bottomrule
\end{cvtable}
}

\newcommand{\resumeSubItem}[2]{\resumeItem{#1}{#2}\vspace{-4pt}}

\renewcommand{\labelitemi}{$\vcenter{\hbox{\tiny$\bullet$}}$}
\renewcommand{\labelitemii}{$\vcenter{\hbox{\tiny$\circ$}}$}

\newcommand{\resumeSubHeadingListStart}{\begin{itemize}[leftmargin=*,labelsep=1mm]}
\newcommand{\resumeHeadingSkillStart}{\begin{itemize}[leftmargin=*,itemsep=1.7mm, rightmargin=2ex]}
\newcommand{\resumeItemListStart}{\begin{itemize}[leftmargin=*,labelsep=1mm,itemsep=0.5mm]}

\newcommand{\resumeSubHeadingListEnd}{\end{itemize}\vspace{2mm}}
\newcommand{\resumeHeadingSkillEnd}{\end{itemize}\vspace{-2mm}}
\newcommand{\resumeItemListEnd}{\end{itemize}\vspace{-2mm}}
\newcommand{\cvsection}[1]{%
\vspace{2mm}
\begin{tcolorbox}
    \textbf{\large #1}
\end{tcolorbox}
    \vspace{-4mm}
}

\newcolumntype{L}{>{\raggedright\arraybackslash}X}%
\newcolumntype{R}{>{\raggedleft\arraybackslash}X}%
\newcolumntype{C}{>{\centering\arraybackslash}X}%

% Commands for icon sizing and positioning
\newcommand{\socialicon}[1]{\raisebox{-0.05em}{\resizebox{!}{1em}{#1}}}
\newcommand{\ieeeicon}[1]{\raisebox{-0.3em}{\resizebox{!}{1.3em}{#1}}}

% Profile icon registry. Known compact account IDs select their service icon;
% arbitrary `{label, url}` links use \cviconlink.
\newcommand{\cviconscholar}{\socialicon{\aiGoogleScholar}}
\newcommand{\cviconorcid}{\socialicon{\aiOrcid}}
\newcommand{\cviconlinkedin}{\socialicon{\faLinkedin}}
\newcommand{\cvicongithub}{\socialicon{\faGithub}}
\newcommand{\cviconlink}{\socialicon{\faLink}}

% Font options
\newcommand{\headerfonti}{\fontfamily{phv}\selectfont} % Helvetica-like (similar to Arial/Calibri)
\newcommand{\headerfontii}{\fontfamily{ptm}\selectfont} % Times-like (similar to Times New Roman)
\newcommand{\headerfontiii}{\rmfamily} % TeX Gyre Pagella, set above with Unicode support
\newcommand{\headerfontiv}{\fontfamily{pbk}\selectfont} % Bookman (readable serif)
\newcommand{\headerfontv}{\fontfamily{pag}\selectfont} % Avant Garde-like (similar to Trebuchet MS)
\newcommand{\headerfontvi}{\fontfamily{cmss}\selectfont} % Computer Modern Sans Serif
\newcommand{\headerfontvii}{\fontfamily{qhv}\selectfont} % Quasi-Helvetica (another Arial/Calibri alternative)
\newcommand{\headerfontviii}{\fontfamily{qpl}\selectfont} % Quasi-Palatino (another elegant serif option)
\newcommand{\headerfontix}{\fontfamily{qtm}\selectfont} % Quasi-Times (another Times New Roman alternative)
\newcommand{\headerfontx}{\fontfamily{bch}\selectfont} % Charter (clean serif font)

% --------------------------------
% BIBLATEX (publications + talks)
% --------------------------------
\usepackage[
  backend=biber,
  style=numeric,
  sorting=ydnt,
  defernumbers=true,
  maxbibnames=99,
  giveninits=true,
  doi=true,
  url=true,
  isbn=false
]{biblatex}

% Two bib databases (kept separate via refsection below).
% Publications are the repository-wide canonical list, shared with the website.
\addbibresource{../content/publications.bib}
\addbibresource{../content/talks.bib}

% The shared bibliography may carry rendering fields used only by the website. They are not part of the CV, and abstracts can contain raw % and &
% that would break the LaTeX pass, so drop them during biber processing instead
% of altering the .bib the website reads.
\DeclareSourcemap{
  \maps[datatype=bibtex]{
    \map{
      \step[fieldset=abstract, null]
      \step[fieldset=selected, null]
      \step[fieldset=preview, null]
      \step[fieldset=abbr, null]
      \step[fieldset=bibtex_show, null]
      \step[fieldset=html, null]
      \step[fieldset=pdf, null]
      \step[fieldset=code, null]
      \step[fieldset=altmetric, null]
      \step[fieldset=dimensions, null]
    }
    % DBLP's @misc artifact records carry an UNFILLED placeholder,
    % `note = {Accessed on YYYY-MM-DD.}`, which biblatex would print verbatim as
    % the venue of a released software package. The website already skips it
    % (`VENUE_FIELDS` in web/src/lib/record.ts prefers `publisher`); this is the
    % same fix on the printed side, so declaring a `types: [misc]` section in
    % content/cv.yaml renders correctly instead of printing the placeholder.
    \map{
      \step[fieldsource=note, match=\regexp{\AAccessed\son\sYYYY-MM-DD\.?\Z}, final]
      \step[fieldset=note, null]
    }
  }
}

% Compact bib
\setlength\bibitemsep{0.25\baselineskip}
\renewcommand*{\bibfont}{\footnotesize}

% Bib headings matching your tidy style
\defbibheading{bibsubheading}[\refname]{%
  \vspace{0.8mm}\textbf{#1}\par\vspace{-0.3mm}
}

% No filters here. Which entry type belongs under which heading is a curated
% opinion, so it is declared in `content/cv.yaml` under `publications:` and
% `talks:` and read by the website too. build-cv-data.mjs turns each declared
% section into its \defbibfilter and its \printbibliography, and gives this file
% \cvPublicationsKey / \cvPublicationsSections and \cvTalksKey /
% \cvTalksSections, plus \cvPublicationsCount / \cvTalksCount - how many entries
% the .bib holds, which guards the \refsection below.
% Adding, reordering or renaming a section is a cv.yaml edit.

% --------------------------------
% CONTENT (generated from content/cv.yaml)
% --------------------------------
% =============================================================================
% GENERATED FILE - DO NOT EDIT.
%
% Produced by `node scripts/build-cv-data.mjs` from `content/cv.yaml`.
% Edit cv.yaml and regenerate; hand edits here are overwritten and CI rejects
% them (the workflow fails if this file is stale relative to cv.yaml).
%
% This file holds CONTENT ONLY. Layout, spacing and styling live in cv/preamble.tex.
% =============================================================================

\newcommand{\cvName}{%
Sahana Aster KŌWHAI, Ph.D.%
}

\newcommand{\cvContactLine}{%
\href{mailto:sahana.kowhai@example.edu}{sahana.kowhai@example.edu}
\;|\;
\href{https://profiles.example.edu/sahana-kowhai}{Staff profile}%
}

\newcommand{\cvProfilesLine}{%
\cviconorcid\,\href{https://orcid.org/0000-0000-0000-0000}{0000-0000-0000-0000}
\;|\;
\cviconscholar\,\href{https://scholar.google.com/citations?user=sahana-kowhai-example}{sahana-kowhai-example}
\;|\;
\cvicongithub\,\href{https://github.com/sahana-kowhai}{sahana-kowhai}
\;|\;
\cviconlinkedin\,\href{https://www.linkedin.com/in/sahana-kowhai}{sahana-kowhai}
\;|\;
\cviconlink\,\href{https://bsky.app/profile/sahana-kowhai.example}{Bluesky}%
}

\newcommand{\cvAffiliation}{%
Associate Professor of Palaeoclimatology, School of Earth and Ice Studies \\
Kārearea University, \\
Tideglass Research Cooperative, Ōtepoti Dunedin, Aotearoa New Zealand%
}

\newcommand{\cvShortBio}{%
Sahana Aster Kōwhai is a \textbf{fictional} palaeoclimatologist whose \emph{research} reconstructs Southern Ocean change from marine sediment archives.%
}

\newcommand{\cvFocus}{%
Southern Ocean palaeoceanography across glacial terminations, with an emphasis on multiproxy reconstructions, reproducible chronology and equitable Antarctic field partnerships.%
}

\newcommand{\cvFooter}{%
\textbf{Example record.} This curriculum vitae describes a fictional scholar and is supplied only to demonstrate the template.%
}

\newcommand{\cvTeachingStatementNote}{%
I teach palaeoclimate reasoning as \textbf{evidence work}: students leave able to defend an age model, not only to run one.%
}

\newcommand{\cvTeachingStatement}{%
%
}

\newcommand{\cvTeachingStatementRows}{%
%
}

\newcommand{\cvTeachingStatementHeader}{%
%
}

\newcommand{\cvTeachingStatementInline}{%
%
}

\newcommand{\cvTeachingStatementCount}{0}

\newcommand{\cvStrandsNote}{%
%
}

\newcommand{\cvStrands}{%
\ifnum1>\cvmax\else
\cventry
  {Ocean archives}
  {Reading marine sediments as layered records of temperature, circulation and sea-ice change.}
  {}
  {}
\resumeItemListStart
  \item Evidence lives in content/publications.bib and the fieldwork section.
\resumeItemListEnd
\fi

\ifnum2>\cvmax\else
\cventry
  {Chronology and uncertainty}
  {Making age models, proxy calibration and uncertainty legible from sample to published claim.}
  {}
  {}
\resumeItemListStart
  \item Evidence lives in content/publications.bib and the teaching record.
\resumeItemListEnd
\fi

\ifnum3>\cvmax\else
\cventry
  {Open polar practice}
  {Designing reusable datasets and reciprocal field partnerships for Southern Ocean research.}
  {}
  {}
\resumeItemListStart
  \item Evidence lives in the projects and service sections.
\resumeItemListEnd
\fi%
}

\newcommand{\cvStrandsRows}{%
Ocean archives & Reading marine sediments as layered records of temperature, circulation and sea-ice change. & Evidence lives in content/publications.bib and the fieldwork section. \\
Chronology and uncertainty & Making age models, proxy calibration and uncertainty legible from sample to published claim. & Evidence lives in content/publications.bib and the teaching record. \\
Open polar practice & Designing reusable datasets and reciprocal field partnerships for Southern Ocean research. & Evidence lives in the projects and service sections. \\%
}

\newcommand{\cvStrandsHeader}{%
\textbf{Title} & \textbf{Detail} & \textbf{Items} \\%
}

\newcommand{\cvStrandsInline}{%
Ocean archives (Reading marine sediments as layered records of temperature, circulation and sea-ice change.), Chronology and uncertainty (Making age models, proxy calibration and uncertainty legible from sample to published claim.), Open polar practice (Designing reusable datasets and reciprocal field partnerships for Southern Ocean research.)%
}

\newcommand{\cvStrandsCount}{3}

\newcommand{\cvStrandsPrinted}{0}

\newcommand{\cvAppointmentsNote}{%
%
}

\newcommand{\cvAppointments}{%
\ifnum1>\cvmax\else
\cventry
  {Associate Professor of Palaeoclimatology}
  {Kārearea University, School of Earth and Ice Studies}
  {Feb 2021 -- Present}
  {Ōtepoti Dunedin, Aotearoa New Zealand}
\resumeItemListStart
  \item Cross-appointed affiliate of the Tideglass Research Cooperative.
\resumeItemListEnd
\fi

\ifnum2>\cvmax\else
\cventry
  {Senior Lecturer in Marine Geology}
  {Aster Coast University, Koru School of Marine Records}
  {Jul 2016 -- Jan 2021}
  {nipaluna Hobart, Australia}
\fi%
}

\newcommand{\cvAppointmentsRows}{%
Associate Professor of Palaeoclimatology & Kārearea University & Ōtepoti Dunedin, Aotearoa New Zealand & Feb 2021 -- Present & School of Earth and Ice Studies & https://earth-ice.example.edu & Cross-appointed affiliate of the Tideglass Research Cooperative. & 2021-02-15 \\
Senior Lecturer in Marine Geology & Aster Coast University & nipaluna Hobart, Australia & Jul 2016 -- Jan 2021 & Koru School of Marine Records & 2016-07-01 \\%
}

\newcommand{\cvAppointmentsHeader}{%
\textbf{Title} & \textbf{Org} & \textbf{Place} & \textbf{Dates} & \textbf{Detail} & \textbf{Url} & \textbf{Items} & \textbf{Announced} \\%
}

\newcommand{\cvAppointmentsInline}{%
Associate Professor of Palaeoclimatology (School of Earth and Ice Studies), Senior Lecturer in Marine Geology (Koru School of Marine Records)%
}

\newcommand{\cvAppointmentsCount}{2}

\newcommand{\cvEducationNote}{%
%
}

\newcommand{\cvEducation}{%
\ifnum1>\cvmax\else
\cventry
  {Ph.D. in Marine Geology}
  {Aster Institute for Earth Records}
  {2010 -- 2014}
  {Canberra, Australia}
\resumeItemListStart
  \item Thesis: Foraminiferal Mg/Ca and Southern Ocean warming across Termination I.
\resumeItemListEnd
\fi

\ifnum2>\cvmax\else
\cventry
  {M.Sc. in Earth Sciences (First Class Honours)}
  {Kōwhai Coast University}
  {2008 -- 2009}
  {Cape Town, South Africa}
\fi

\ifnum3>\cvmax\else
\cventry
  {B.Sc. in Geology and Oceanography}
  {Tideglass Ocean Institute}
  {2005 -- 2008}
  {Chennai, India}
\fi%
}

\newcommand{\cvEducationRows}{%
Ph.D. in Marine Geology & Aster Institute for Earth Records & Canberra, Australia & 2010 -- 2014 & Thesis: Foraminiferal Mg/Ca and Southern Ocean warming across Termination I. \\
M.Sc. in Earth Sciences (First Class Honours) & Kōwhai Coast University & Cape Town, South Africa & 2008 -- 2009 \\
B.Sc. in Geology and Oceanography & Tideglass Ocean Institute & Chennai, India & 2005 -- 2008 \\%
}

\newcommand{\cvEducationHeader}{%
\textbf{Title} & \textbf{Org} & \textbf{Place} & \textbf{Dates} & \textbf{Items} \\%
}

\newcommand{\cvEducationInline}{%
Ph.D. in Marine Geology, M.Sc. in Earth Sciences (First Class Honours), B.Sc. in Geology and Oceanography%
}

\newcommand{\cvEducationCount}{3}

\newcommand{\cvTeachingNote}{%
%
}

\newcommand{\cvTeaching}{%
\ifnum1>\cvmax\else
\cventry
  {Course coordinator and lecturer}
  {Kārearea University}
  {2021 -- Present}
  {Ōtepoti Dunedin, Aotearoa New Zealand}
\cvcourses{Y Y Y Y}{\textbf{Course} & \textbf{Programme} & \textbf{Topics} & \textbf{Credit points} \\}{
GEOL 273 & B.Sc. Geology & Palaeoclimate archives; proxy calibration & 18 \\
GEOL 472 & B.Sc. (Hons) & Marine cores; age-depth modelling & 18 \\
OCEN 405 & M.Sc. Oceanography & Southern Ocean circulation; uncertainty & 20 \\}
\fi%
}

\newcommand{\cvTeachingRows}{%
Course coordinator and lecturer & Kārearea University & Ōtepoti Dunedin, Aotearoa New Zealand & 2021 -- Present & GEOL 273 B.Sc. Geology Palaeoclimate archives; proxy calibration 18, GEOL 472 B.Sc. (Hons) Marine cores; age-depth modelling 18, OCEN 405 M.Sc. Oceanography Southern Ocean circulation; uncertainty 20 \\%
}

\newcommand{\cvTeachingHeader}{%
\textbf{Title} & \textbf{Org} & \textbf{Place} & \textbf{Dates} & \textbf{Rows} \\%
}

\newcommand{\cvTeachingInline}{%
Course coordinator and lecturer%
}

\newcommand{\cvTeachingCount}{1}

\newcommand{\cvSupervisionNote}{%
\textbf{Current supervision:} projects span marine geochemistry, chronology and open polar data.%
}

\newcommand{\cvSupervision}{%
\ifnum1>\cvmax\else
\cventry
  {Doctoral candidates}
  {Primary supervisor for two; co-supervisor for two, 4}
  {}
  {}
\fi

\ifnum2>\cvmax\else
\cventry
  {M.Sc. and honours researchers}
  {Completed and current research dissertations, 9}
  {}
  {}
\fi%
}

\newcommand{\cvSupervisionRows}{%
Doctoral candidates & 4 & Primary supervisor for two; co-supervisor for two \\
M.Sc. and honours researchers & 9 & Completed and current research dissertations \\%
}

\newcommand{\cvSupervisionHeader}{%
\textbf{Title} & \textbf{Count} & \textbf{Detail} \\%
}

\newcommand{\cvSupervisionInline}{%
Doctoral candidates (Primary supervisor for two; co-supervisor for two), M.Sc. and honours researchers (Completed and current research dissertations)%
}

\newcommand{\cvSupervisionCount}{2}

\newcommand{\cvServiceNote}{%
%
}

\newcommand{\cvService}{%
\ifnum1>\cvmax\else
\cventry
  {Associate Editor}
  {Tideglass Climate Letters, Marine palaeoclimate}
  {2023 -- Present}
  {}
\fi

\ifnum2>\cvmax\else
\cventry
  {Associate Editor}
  {Example Journal of Palaeoclimate Studies, 2022--2024 \textbf{[IF: 3.4, Q1]}}
  {}
  {}
\fi

\ifnum3>\cvmax\else
\cventry
  {Scientific Steering Committee}
  {Tideglass Research Cooperative, 2022--2026}
  {}
  {}
\fi

\ifnum4>\cvmax\else
\cventry
  {Panel reviewer}
  {Kōwhai Research Trust}
  {2021 -- Present}
  {}
\fi

\ifnum5>\cvmax\else
\cventry
  {Reviewer}
  {Open Core Review}
  {}
  {}
\fi%
}

\newcommand{\cvServiceRows}{%
Associate Editor & Tideglass Climate Letters & 2023 -- Present & Marine palaeoclimate & 2023-03-10 \\
Associate Editor & Example Journal of Palaeoclimate Studies & 2022--2024 & IF: 3.4, Q1 & https://example.org/journals/palaeoclimate-studies \\
Scientific Steering Committee & Tideglass Research Cooperative & 2022--2026 \\
Panel reviewer & Kōwhai Research Trust & 2021 -- Present \\
Reviewer & Open Core Review \\%
}

\newcommand{\cvServiceHeader}{%
\textbf{Title} & \textbf{Org} & \textbf{Dates} & \textbf{Detail} & \textbf{Announced} \\%
}

\newcommand{\cvServiceInline}{%
Associate Editor (Marine palaeoclimate), Associate Editor, Scientific Steering Committee, Panel reviewer, Reviewer%
}

\newcommand{\cvServiceCount}{5}

\newcommand{\cvProjectsNote}{%
%
}

\newcommand{\cvProjects}{%
\ifnum1>\cvmax\else
\cventry
  {TIDE --- Termination Ice Dynamics and Ecosystems}
  {Kōwhai Research Trust, Principal investigator}
  {2024 -- 2027}
  {}
\resumeItemListStart
  \item Multiproxy reconstruction of Southern Ocean fronts across Termination II, with voyage support from the \href{https://voyages.example.edu}{Southern Archive Voyage Consortium}.
\resumeItemListEnd
\fi

\ifnum2>\cvmax\else
\cventry
  {Open Core Chronologies}
  {Koru Basin Research Alliance, Co-investigator}
  {2022 -- 2025}
  {}
\resumeItemListStart
  \item Reusable age-depth models and provenance for legacy sediment cores.
\resumeItemListEnd
\fi%
}

\newcommand{\cvProjectsRows}{%
TIDE --- Termination Ice Dynamics and Ecosystems & Kōwhai Research Trust & 2024 -- 2027 & Principal investigator & NZ\$960,000 & https://research.example.edu/tide & 2024-06-03 & Multiproxy reconstruction of Southern Ocean fronts across Termination II, with voyage support from the \href{https://voyages.example.edu}{Southern Archive Voyage Consortium}. \\
Open Core Chronologies & Koru Basin Research Alliance & 2022 -- 2025 & Co-investigator & A\$420,000 & Reusable age-depth models and provenance for legacy sediment cores. \\%
}

\newcommand{\cvProjectsHeader}{%
\textbf{Title} & \textbf{Org} & \textbf{Dates} & \textbf{Detail} & \textbf{Funding} & \textbf{Url} & \textbf{Announced} & \textbf{Items} \\%
}

\newcommand{\cvProjectsInline}{%
TIDE --- Termination Ice Dynamics and Ecosystems (Principal investigator), Open Core Chronologies (Co-investigator)%
}

\newcommand{\cvProjectsCount}{2}

\newcommand{\cvFieldworkNote}{%
%
}

\newcommand{\cvFieldwork}{%
\ifnum1>\cvmax\else
\cventry
  {Pukaki Rise sediment coring voyage}
  {RV Tāwhirimātea, Co-chief scientist}
  {Jan -- Feb 2025}
  {Southern Ocean}
\fi

\ifnum2>\cvmax\else
\cventry
  {East Antarctic margin core curation}
  {Koru Basin Core Archive, Visiting scientist}
  {Nov 2022}
  {Te Whanganui-a-Tara Wellington}
\fi%
}

\newcommand{\cvFieldworkRows}{%
Pukaki Rise sediment coring voyage & RV Tāwhirimātea & Southern Ocean & Jan -- Feb 2025 & Co-chief scientist & 2025-01-12 \\
East Antarctic margin core curation & Koru Basin Core Archive & Te Whanganui-a-Tara Wellington & Nov 2022 & Visiting scientist \\%
}

\newcommand{\cvFieldworkHeader}{%
\textbf{Title} & \textbf{Org} & \textbf{Place} & \textbf{Dates} & \textbf{Detail} & \textbf{Announced} \\%
}

\newcommand{\cvFieldworkInline}{%
Pukaki Rise sediment coring voyage (Co-chief scientist), East Antarctic margin core curation (Visiting scientist)%
}

\newcommand{\cvFieldworkCount}{2}

\newcommand{\cvFieldworkPrinted}{0}

\newcommand{\cvOutreachNote}{%
%
}

\newcommand{\cvOutreach}{%
\ifnum1>\cvmax\else
\cventry
  {Sediment stories: reading climate from ocean cores}
  {Kōwhai Coast Learning Collective, Community workshop for secondary-school science teachers}
  {Sep 2025}
  {Ōamaru, Aotearoa New Zealand}
\fi

\ifnum2>\cvmax\else
\cventry
  {Core archive open day}
  {Aotearoa Marine Memory Trust, Public demonstration of core scanning and age-model uncertainty}
  {Mar 2024}
  {Porirua, Aotearoa New Zealand}
\fi

\ifnum3>\cvmax\else
\cventry
  {Antarctic archives teacher residency}
  {Polar Evidence Studio, Short course on proxy evidence for secondary-school teachers}
  {Jul 2023}
  {Ōtepoti Dunedin, Aotearoa New Zealand}
\fi%
}

\newcommand{\cvOutreachRows}{%
Sediment stories: reading climate from ocean cores & Kōwhai Coast Learning Collective & Ōamaru, Aotearoa New Zealand & Sep 2025 & Community workshop for secondary-school science teachers & https://learning.example.edu/sediment-stories \\
Core archive open day & Aotearoa Marine Memory Trust & Porirua, Aotearoa New Zealand & Mar 2024 & Public demonstration of core scanning and age-model uncertainty \\
Antarctic archives teacher residency & Polar Evidence Studio & Ōtepoti Dunedin, Aotearoa New Zealand & Jul 2023 & Short course on proxy evidence for secondary-school teachers \\%
}

\newcommand{\cvOutreachHeader}{%
\textbf{Title} & \textbf{Org} & \textbf{Place} & \textbf{Dates} & \textbf{Detail} & \textbf{Url} \\%
}

\newcommand{\cvOutreachInline}{%
Sediment stories: reading climate from ocean cores (Community workshop for secondary-school science teachers), Core archive open day (Public demonstration of core scanning and age-model uncertainty), Antarctic archives teacher residency (Short course on proxy evidence for secondary-school teachers)%
}

\newcommand{\cvOutreachCount}{3}

\newcommand{\cvAwardsNote}{%
%
}

\newcommand{\cvAwards}{%
\ifnum1>\cvmax\else
\cventry
  {Early Career Research Excellence Award}
  {Kārearea Earth Forum}
  {2019}
  {}
\fi

\ifnum2>\cvmax\else
\cventry
  {Tideglass Visiting Fellowship}
  {Tideglass Research Cooperative}
  {2017}
  {Cambridge, United Kingdom}
\fi%
}

\newcommand{\cvAwardsRows}{%
Early Career Research Excellence Award & Kārearea Earth Forum & 2019 & 2019-11-20 \\
Tideglass Visiting Fellowship & Tideglass Research Cooperative & Cambridge, United Kingdom & 2017 \\%
}

\newcommand{\cvAwardsHeader}{%
\textbf{Title} & \textbf{Org} & \textbf{Dates} & \textbf{Announced} \\%
}

\newcommand{\cvAwardsInline}{%
Early Career Research Excellence Award, Tideglass Visiting Fellowship%
}

\newcommand{\cvAwardsCount}{2}

\newcommand{\cvLeadershipNote}{%
%
}

\newcommand{\cvLeadership}{%
\ifnum1>\cvmax\else
\cventry
  {Co-chair, Open Polar Data Working Group}
  {Tideglass Research Cooperative}
  {2024 -- Present}
  {}
\fi

\ifnum2>\cvmax\else
\cventry
  {Faculty representative --- \href{/blog/2026/reading-a-core/}{field note}}
  {Pacific research partnerships committee}
  {2022 -- 2024}
  {}
\fi%
}

\newcommand{\cvLeadershipRows}{%
Co-chair, Open Polar Data Working Group & Tideglass Research Cooperative & 2024 -- Present & 2024-02-01 \\
Faculty representative --- \href{/blog/2026/reading-a-core/}{field note} & Pacific research partnerships committee & 2022 -- 2024 \\%
}

\newcommand{\cvLeadershipHeader}{%
\textbf{Title} & \textbf{Org} & \textbf{Dates} & \textbf{Announced} \\%
}

\newcommand{\cvLeadershipInline}{%
Co-chair, Open Polar Data Working Group, Faculty representative --- \href{/blog/2026/reading-a-core/}{field note}%
}

\newcommand{\cvLeadershipCount}{2}

\newcommand{\cvLeadershipPrinted}{0}

\newcommand{\cvLanguagesNote}{%
%
}

\newcommand{\cvLanguages}{%
\ifnum1>\cvmax\else
\cventry
  {Tamil}
  {Native}
  {}
  {}
\fi

\ifnum2>\cvmax\else
\cventry
  {English}
  {Fluent}
  {}
  {}
\fi

\ifnum3>\cvmax\else
\cventry
  {te reo Māori}
  {Intermediate learner}
  {}
  {}
\fi%
}

\newcommand{\cvLanguagesRows}{%
Tamil & Native \\
English & Fluent \\
te reo Māori & Intermediate learner \\%
}

\newcommand{\cvLanguagesHeader}{%
\textbf{Title} & \textbf{Detail} \\%
}

\newcommand{\cvLanguagesInline}{%
Tamil (Native), English (Fluent), te reo Māori (Intermediate learner)%
}

\newcommand{\cvLanguagesCount}{3}

\newcommand{\cvPublicationsCount}{10}

\defbibfilter{Publications1}{type=article and not keyword=underreview}
\defbibfilter{Publications2}{type=inproceedings and not ( type=article and not keyword=underreview )}
\defbibfilter{Publications3}{( type=incollection or type=inbook or type=book ) and not ( type=article and not keyword=underreview ) and not ( type=inproceedings )}
\defbibfilter{Publications4}{type=dataset and not ( type=article and not keyword=underreview ) and not ( type=inproceedings ) and not ( ( type=incollection or type=inbook or type=book ) )}
\defbibfilter{Publications6}{type=article and keyword=underreview and not ( type=article and not keyword=underreview ) and not ( type=inproceedings ) and not ( ( type=incollection or type=inbook or type=book ) ) and not ( type=dataset ) and not ( type=misc )}

\newcommand{\cvPublicationsKey}{%
J=Journal, C=Conference, B=Books, D=Data, R=Under review%
}

\newcommand{\cvPublicationsSections}{%
\begin{refcontext}[labelprefix=J]
\printbibliography[heading=bibsubheading, title={Journal articles (peer-reviewed)}, filter=Publications1, resetnumbers=true]
\end{refcontext}

\begin{refcontext}[labelprefix=C]
\printbibliography[heading=bibsubheading, title={Conference papers (peer-reviewed)}, filter=Publications2, resetnumbers=true]
\end{refcontext}

\begin{refcontext}[labelprefix=B]
\printbibliography[heading=bibsubheading, title={Books and chapters}, filter=Publications3, resetnumbers=true]
\end{refcontext}

\begin{refcontext}[labelprefix=D]
\printbibliography[heading=bibsubheading, title={Data archives}, filter=Publications4, resetnumbers=true]
\end{refcontext}

\begin{refcontext}[labelprefix=R]
\printbibliography[heading=bibsubheading, title={Manuscripts under review}, filter=Publications6, resetnumbers=true]
\end{refcontext}%
}

\newcommand{\cvTalksCount}{4}

\defbibfilter{Talks1}{type=unpublished and keyword=invited}

\newcommand{\cvTalksKey}{%
I=Invited%
}

\newcommand{\cvTalksSections}{%
\begin{refcontext}[labelprefix=I]
\printbibliography[heading=bibsubheading, title={Invited and keynote talks}, filter=Talks1, resetnumbers=true]
\end{refcontext}%
}

\newcommand{\cvAutoSections}{%
\cvautopart{Teaching Statement}{TeachingStatement}%
\cvautopart{Strands}{Strands}%
\cvautopart{Appointments}{Appointments}%
\cvautopart{Education}{Education}%
\cvautopart{Teaching}{Teaching}%
\cvautopart{Supervision}{Supervision}%
\cvautopart{Service}{Service}%
\cvautopart{Projects}{Projects}%
\cvautopart{Fieldwork}{Fieldwork}%
\cvautopart{Public engagement}{Outreach}%
\cvautopart{Awards}{Awards}%
\cvautopart{Leadership}{Leadership}%
\cvautopart{Languages}{Languages}%
}


,
% so a change to the printed design is made once and reaches all of them.
%
% To change what a CV says, edit content/cv.yaml and regenerate. To change how
% it looks, edit this file. To change WHICH of it one document prints, edit that
% document's body - cv/cv.tex, cv/short.tex or cv/teaching.tex.
%
% WHICH sections print is a record fact, not a layout one. The `\cvpart` lines in
% a document body places the sections whose heading or setting that document
% curates; `\cvAutoSections` then prints every remaining section of
% content/cv.yaml, in the order the record writes them, under the heading its key
% spells out. Adding a section to the record therefore needs no edit here. Give
% a section `printed: false` in content/cv.yaml to keep it off the PDF - that
% holds for the curated ones below too - and `heading:` to name an uncurated one
% something other than its key.
%
% Build:  latexmk -xelatex -cd cv/cv.tex      (xelatex is required: academicons)
%         the same for cv/short.tex and cv/teaching.tex
% -----------------------------------------------------------------------------



\documentclass[a4paper,11pt]{article}

% Package imports
\usepackage{latexsym}
\usepackage{xcolor}
\usepackage{float}
\usepackage{ragged2e}
\usepackage[empty]{fullpage}
\usepackage{wrapfig}
\usepackage{lipsum}
\usepackage{booktabs}   % high-quality horizontal rules (no vertical rules)
\usepackage{array}      % better column types
\usepackage{xparse}
\usepackage{tabularx}
\usepackage{titlesec}
\usepackage{geometry}
\usepackage{marvosym}
\usepackage{verbatim}
\usepackage{enumitem}
\usepackage{fancyhdr}
\usepackage{multicol}
\usepackage{graphicx}
\usepackage{fontspec}
\setmainfont{TeX Gyre Pagella}
\usepackage{fontawesome5}
\usepackage{academicons} % <-- NEW (Scholar, ORCID, etc.)
\usepackage{eurosym} % <-- NEW (€)

% Column helpers (tight but readable)
\newcolumntype{P}[1]{>{\raggedright\arraybackslash}p{#1}}
\newcolumntype{Y}{>{\raggedright\arraybackslash}X}

% A reusable CV table environment (no floats, correct spacing)
\NewDocumentEnvironment{cvtable}{m +b}{%
  \vspace{0.25em}
  \small
  \setlength{\tabcolsep}{5pt}
  \renewcommand{\arraystretch}{1.15}
  \begin{tabularx}{\linewidth}{@{}#1@{}}
  #2
  \end{tabularx}
  \vspace{0.25em}
}{}

% Color definitions
\definecolor{darkblue}{RGB}{0,0,139}

% Page layout
\setlength{\multicolsep}{0pt}
\pagestyle{fancy}
\fancyhf{} % clear all header and footer fields
\fancyfoot{}
\renewcommand{\headrulewidth}{0pt}
\renewcommand{\footrulewidth}{0pt}
\geometry{left=1.4cm, top=0.8cm, right=1.2cm, bottom=1cm}
\setlength{\footskip}{5pt} % Addressing fancyhdr warning

% Hyperlink setup (moved after fancyhdr to address warning)
\usepackage[hidelinks]{hyperref}
\hypersetup{
    colorlinks=true,
    linkcolor=darkblue,
    filecolor=darkblue,
    urlcolor=darkblue,
}

% Custom box settings
\usepackage[most]{tcolorbox}
\tcbset{
    frame code={},
    center title,
    left=0pt,
    right=0pt,
    top=0pt,
    bottom=0pt,
    colback=gray!20,
    colframe=white,
    width=\dimexpr\textwidth\relax,
    enlarge left by=-2mm,
    boxsep=4pt,
    arc=0pt,outer arc=0pt,
}

% URL style
\urlstyle{same}

% Text alignment
\raggedright
\setlength{\tabcolsep}{0in}

% Section formatting
\titleformat{\section}{
  \vspace{-4pt}\scshape\raggedright\large
}{}{0em}{}[\color{black}\titlerule \vspace{-7pt}]

% Custom commands
\newcommand{\resumeItem}[2]{
  \item{
    \textbf{#1}{\hspace{0.5mm}#2 \vspace{-0.5mm}}
  }
}

\newcommand{\resumePOR}[3]{
\vspace{0.5mm}\item
    \begin{tabular*}{0.97\textwidth}[t]{l@{\extracolsep{\fill}}r}
        \textbf{#1}\hspace{0.3mm}#2 & \textit{\small{#3}}
    \end{tabular*}
    \vspace{-2mm}
}

% One entry, every section: what it was, where it was, when, and the place.
% `content/cv.yaml` has exactly one entry shape and this is how it is set.
\newcommand{\cventry}[4]{
\vspace{0.5mm}\item
    \begin{tabularx}{0.98\textwidth}[t]{@{}Y >{\raggedleft\arraybackslash}p{4.8cm}@{}}
        \textbf{#1} & \textit{\footnotesize{#3}} \\
        \textit{\footnotesize{#2}} &  \footnotesize{#4}\\
    \end{tabularx}
    \vspace{-2.4mm}
}

% An empty generated macro, to test the optional prose fields against.
% NOT `\ifx\cvFocus\empty`: \newcommand makes its macros \long and \empty is
% not, so \ifx compares false for two macros that both hold nothing, and the
% heading prints over an empty field. This one is made the same way they are.
\newcommand{\cvempty}{}

% Separates the paragraphs of a section's `note:`. Spacing is layout, so the
% amount of it lives here and the generator only says where the break goes.
\newcommand{\cvnotesep}{\par\vspace{1.2mm}}

% Record that this file lays out one section key itself, so \cvautopart below
% does not print it a second time from the record's own section sequence.
% \gdef, not \def: \cvpartflush sets its entries inside a group.
\newcommand{\cvlaidout}[1]{\expandafter\gdef\csname cvlaidout@#1\endcsname{}}

% Make the five generated macros this layout may mention safe even when
% content/cv.yaml has no such section, and default the section to printing.
% (`Header` is the sixth generated macro, but this layout supplies its own table
% headings.) `\cvpart` does this for you; a hand-written section block must call it
% itself.
\newcommand{\cvdeclare}[1]{%
  \cvlaidout{#1}%
  % Printing is the default: only `printed: false` in content/cv.yaml generates a
  % \cv<Key>Printed of 0, and this \providecommand leaves that definition alone.
  \expandafter\providecommand\csname cv#1Printed\endcsname{1}%
  \expandafter\providecommand\csname cv#1Count\endcsname{0}%
  \expandafter\providecommand\csname cv#1\endcsname{}%
  \expandafter\providecommand\csname cv#1Note\endcsname{}%
  \expandafter\providecommand\csname cv#1Rows\endcsname{}%
  \expandafter\providecommand\csname cv#1Inline\endcsname{}%
}

% The same, for a bibliography section: `publications:`/`talks:` in cv.yaml
% generate \cv<Name>Key, \cv<Name>Sections and \cv<Name>Count (the number of
% entries in the matching .bib) instead of the five above. An adopter who has
% declared neither - which is everyone on day one - would otherwise meet a raw
% "Undefined control sequence" that stops xelatex at its interactive prompt.
% Empty is the documented behaviour: the heading prints over an empty
% bibliography and README.md tells them to delete the block.
\newcommand{\cvdeclarebib}[1]{%
  \expandafter\providecommand\csname cv#1Count\endcsname{0}%
  \expandafter\providecommand\csname cv#1Key\endcsname{}%
  \expandafter\providecommand\csname cv#1Sections\endcsname{}%
}

% How many entries of a section a document prints. The generated \cv<Name>
% guards each of its entries with \ifnum<n>>\cvmax, so this is the cutoff.
%
% This is the layout's call, and it is a different question from the record's
% `printed:`. WHICH sections print is a record fact; HOW MANY entries of one a
% document prints is a page budget belonging to that document - the same section
% is worth three entries in a short CV and all of them in the full one - so the
% number lives here and content/cv.yaml never states it. A count high enough
% that no record reaches it is the default, which keeps the full CV whole.
%
% Curation is neither: "my five best papers" is a judgement about the work, so
% it is marked on the fact itself - see the `selected` keyword in
% content/publications.bib and how cv/short.tex filters on it.
\newcommand{\cvmax}{99999}

% Print one section of content/cv.yaml. #2 is the heading, #3 the section's
% macro name (`Awards` for the `awards:` key), and the optional #1 is how many
% of its entries to print. A section the file does not have, or gives
% `printed: false`, prints nothing at all - no heading, no gap - so this line
% can stay. Reorder the CV by reordering these lines in the document body.
%
% The whole thing is one group so that \cvmax is restored afterwards: a
% \cvpart[3] must not silently truncate the section written after it.
\newcommand{\cvpart}[3][99999]{{%
  \renewcommand{\cvmax}{#1}%
  \cvdeclare{#3}%
  \ifnum\csname cv#3Count\endcsname>0
  \ifnum\csname cv#3Printed\endcsname>0
  \section{\textbf{#2}}
  \vspace{-0.4mm}
  \resumeSubHeadingListStart
  \csname cv#3\endcsname
  \resumeSubHeadingListEnd
  \vspace{-6mm}
  \fi
  \fi
}}

% One section of the record's own section sequence: `\cvpart`, unless this file
% has already laid that key out by hand above. Every line of \cvAutoSections is
% one of these, so an adopter's new section prints with no edit here, and the
% ones curated below keep their curated heading, order and setting.
\newcommand{\cvautopart}[2]{%
  \expandafter\ifx\csname cvlaidout@#2\endcsname\relax\cvpart{#1}{#2}\fi%
}

% `\cvpart`, with the second line of every entry set the other way. Which
% sections use which is a layout choice and lives here; cv.yaml says nothing
% about it.
\newcommand{\cvpartflush}[3][99999]{{%
  \renewcommand{\cventry}[4]{%
    \vspace{0.5mm}\item
    \begin{tabularx}{0.98\textwidth}[t]{@{}Y >{\raggedleft\arraybackslash}p{4.8cm}@{}}
        \textbf{##1} & \textit{\footnotesize{##3}} \\
        \footnotesize{\textit{##2}} & \footnotesize{##4}
    \end{tabularx}
    \vspace{-2.4mm}
  }%
  \cvpart[#1]{#2}{#3}}}

% A table hanging under an entry: `rows:` in cv.yaml. #1 is a generated
% equal-width column specification, #2 the generated headings and #3 the rows.
% Any number of consistently shaped columns therefore works without editing
% this layout. The generator refuses rows whose key order disagrees.
\newcommand{\cvcourses}[3]{%
\par
\begin{cvtable}{#1}
\toprule
#2
\midrule
#3
\bottomrule
\end{cvtable}
}

\newcommand{\resumeSubItem}[2]{\resumeItem{#1}{#2}\vspace{-4pt}}

\renewcommand{\labelitemi}{$\vcenter{\hbox{\tiny$\bullet$}}$}
\renewcommand{\labelitemii}{$\vcenter{\hbox{\tiny$\circ$}}$}

\newcommand{\resumeSubHeadingListStart}{\begin{itemize}[leftmargin=*,labelsep=1mm]}
\newcommand{\resumeHeadingSkillStart}{\begin{itemize}[leftmargin=*,itemsep=1.7mm, rightmargin=2ex]}
\newcommand{\resumeItemListStart}{\begin{itemize}[leftmargin=*,labelsep=1mm,itemsep=0.5mm]}

\newcommand{\resumeSubHeadingListEnd}{\end{itemize}\vspace{2mm}}
\newcommand{\resumeHeadingSkillEnd}{\end{itemize}\vspace{-2mm}}
\newcommand{\resumeItemListEnd}{\end{itemize}\vspace{-2mm}}
\newcommand{\cvsection}[1]{%
\vspace{2mm}
\begin{tcolorbox}
    \textbf{\large #1}
\end{tcolorbox}
    \vspace{-4mm}
}

\newcolumntype{L}{>{\raggedright\arraybackslash}X}%
\newcolumntype{R}{>{\raggedleft\arraybackslash}X}%
\newcolumntype{C}{>{\centering\arraybackslash}X}%

% Commands for icon sizing and positioning
\newcommand{\socialicon}[1]{\raisebox{-0.05em}{\resizebox{!}{1em}{#1}}}
\newcommand{\ieeeicon}[1]{\raisebox{-0.3em}{\resizebox{!}{1.3em}{#1}}}

% Profile icon registry. Known compact account IDs select their service icon;
% arbitrary `{label, url}` links use \cviconlink.
\newcommand{\cviconscholar}{\socialicon{\aiGoogleScholar}}
\newcommand{\cviconorcid}{\socialicon{\aiOrcid}}
\newcommand{\cviconlinkedin}{\socialicon{\faLinkedin}}
\newcommand{\cvicongithub}{\socialicon{\faGithub}}
\newcommand{\cviconlink}{\socialicon{\faLink}}

% Font options
\newcommand{\headerfonti}{\fontfamily{phv}\selectfont} % Helvetica-like (similar to Arial/Calibri)
\newcommand{\headerfontii}{\fontfamily{ptm}\selectfont} % Times-like (similar to Times New Roman)
\newcommand{\headerfontiii}{\rmfamily} % TeX Gyre Pagella, set above with Unicode support
\newcommand{\headerfontiv}{\fontfamily{pbk}\selectfont} % Bookman (readable serif)
\newcommand{\headerfontv}{\fontfamily{pag}\selectfont} % Avant Garde-like (similar to Trebuchet MS)
\newcommand{\headerfontvi}{\fontfamily{cmss}\selectfont} % Computer Modern Sans Serif
\newcommand{\headerfontvii}{\fontfamily{qhv}\selectfont} % Quasi-Helvetica (another Arial/Calibri alternative)
\newcommand{\headerfontviii}{\fontfamily{qpl}\selectfont} % Quasi-Palatino (another elegant serif option)
\newcommand{\headerfontix}{\fontfamily{qtm}\selectfont} % Quasi-Times (another Times New Roman alternative)
\newcommand{\headerfontx}{\fontfamily{bch}\selectfont} % Charter (clean serif font)

% --------------------------------
% BIBLATEX (publications + talks)
% --------------------------------
\usepackage[
  backend=biber,
  style=numeric,
  sorting=ydnt,
  defernumbers=true,
  maxbibnames=99,
  giveninits=true,
  doi=true,
  url=true,
  isbn=false
]{biblatex}

% Two bib databases (kept separate via refsection below).
% Publications are the repository-wide canonical list, shared with the website.
\addbibresource{../content/publications.bib}
\addbibresource{../content/talks.bib}

% The shared bibliography may carry rendering fields used only by the website. They are not part of the CV, and abstracts can contain raw % and &
% that would break the LaTeX pass, so drop them during biber processing instead
% of altering the .bib the website reads.
\DeclareSourcemap{
  \maps[datatype=bibtex]{
    \map{
      \step[fieldset=abstract, null]
      \step[fieldset=selected, null]
      \step[fieldset=preview, null]
      \step[fieldset=abbr, null]
      \step[fieldset=bibtex_show, null]
      \step[fieldset=html, null]
      \step[fieldset=pdf, null]
      \step[fieldset=code, null]
      \step[fieldset=altmetric, null]
      \step[fieldset=dimensions, null]
    }
    % DBLP's @misc artifact records carry an UNFILLED placeholder,
    % `note = {Accessed on YYYY-MM-DD.}`, which biblatex would print verbatim as
    % the venue of a released software package. The website already skips it
    % (`VENUE_FIELDS` in web/src/lib/record.ts prefers `publisher`); this is the
    % same fix on the printed side, so declaring a `types: [misc]` section in
    % content/cv.yaml renders correctly instead of printing the placeholder.
    \map{
      \step[fieldsource=note, match=\regexp{\AAccessed\son\sYYYY-MM-DD\.?\Z}, final]
      \step[fieldset=note, null]
    }
  }
}

% Compact bib
\setlength\bibitemsep{0.25\baselineskip}
\renewcommand*{\bibfont}{\footnotesize}

% Bib headings matching your tidy style
\defbibheading{bibsubheading}[\refname]{%
  \vspace{0.8mm}\textbf{#1}\par\vspace{-0.3mm}
}

% No filters here. Which entry type belongs under which heading is a curated
% opinion, so it is declared in `content/cv.yaml` under `publications:` and
% `talks:` and read by the website too. build-cv-data.mjs turns each declared
% section into its \defbibfilter and its \printbibliography, and gives this file
% \cvPublicationsKey / \cvPublicationsSections and \cvTalksKey /
% \cvTalksSections, plus \cvPublicationsCount / \cvTalksCount - how many entries
% the .bib holds, which guards the \refsection below.
% Adding, reordering or renaming a section is a cv.yaml edit.

% --------------------------------
% CONTENT (generated from content/cv.yaml)
% --------------------------------
% =============================================================================
% GENERATED FILE - DO NOT EDIT.
%
% Produced by `node scripts/build-cv-data.mjs` from `content/cv.yaml`.
% Edit cv.yaml and regenerate; hand edits here are overwritten and CI rejects
% them (the workflow fails if this file is stale relative to cv.yaml).
%
% This file holds CONTENT ONLY. Layout, spacing and styling live in cv/preamble.tex.
% =============================================================================

\newcommand{\cvName}{%
Sahana Aster KŌWHAI, Ph.D.%
}

\newcommand{\cvContactLine}{%
\href{mailto:sahana.kowhai@example.edu}{sahana.kowhai@example.edu}
\;|\;
\href{https://profiles.example.edu/sahana-kowhai}{Staff profile}%
}

\newcommand{\cvProfilesLine}{%
\cviconorcid\,\href{https://orcid.org/0000-0000-0000-0000}{0000-0000-0000-0000}
\;|\;
\cviconscholar\,\href{https://scholar.google.com/citations?user=sahana-kowhai-example}{sahana-kowhai-example}
\;|\;
\cvicongithub\,\href{https://github.com/sahana-kowhai}{sahana-kowhai}
\;|\;
\cviconlinkedin\,\href{https://www.linkedin.com/in/sahana-kowhai}{sahana-kowhai}
\;|\;
\cviconlink\,\href{https://bsky.app/profile/sahana-kowhai.example}{Bluesky}%
}

\newcommand{\cvAffiliation}{%
Associate Professor of Palaeoclimatology, School of Earth and Ice Studies \\
Kārearea University, \\
Tideglass Research Cooperative, Ōtepoti Dunedin, Aotearoa New Zealand%
}

\newcommand{\cvShortBio}{%
Sahana Aster Kōwhai is a \textbf{fictional} palaeoclimatologist whose \emph{research} reconstructs Southern Ocean change from marine sediment archives.%
}

\newcommand{\cvFocus}{%
Southern Ocean palaeoceanography across glacial terminations, with an emphasis on multiproxy reconstructions, reproducible chronology and equitable Antarctic field partnerships.%
}

\newcommand{\cvFooter}{%
\textbf{Example record.} This curriculum vitae describes a fictional scholar and is supplied only to demonstrate the template.%
}

\newcommand{\cvTeachingStatementNote}{%
I teach palaeoclimate reasoning as \textbf{evidence work}: students leave able to defend an age model, not only to run one.%
}

\newcommand{\cvTeachingStatement}{%
%
}

\newcommand{\cvTeachingStatementRows}{%
%
}

\newcommand{\cvTeachingStatementHeader}{%
%
}

\newcommand{\cvTeachingStatementInline}{%
%
}

\newcommand{\cvTeachingStatementCount}{0}

\newcommand{\cvStrandsNote}{%
%
}

\newcommand{\cvStrands}{%
\ifnum1>\cvmax\else
\cventry
  {Ocean archives}
  {Reading marine sediments as layered records of temperature, circulation and sea-ice change.}
  {}
  {}
\resumeItemListStart
  \item Evidence lives in content/publications.bib and the fieldwork section.
\resumeItemListEnd
\fi

\ifnum2>\cvmax\else
\cventry
  {Chronology and uncertainty}
  {Making age models, proxy calibration and uncertainty legible from sample to published claim.}
  {}
  {}
\resumeItemListStart
  \item Evidence lives in content/publications.bib and the teaching record.
\resumeItemListEnd
\fi

\ifnum3>\cvmax\else
\cventry
  {Open polar practice}
  {Designing reusable datasets and reciprocal field partnerships for Southern Ocean research.}
  {}
  {}
\resumeItemListStart
  \item Evidence lives in the projects and service sections.
\resumeItemListEnd
\fi%
}

\newcommand{\cvStrandsRows}{%
Ocean archives & Reading marine sediments as layered records of temperature, circulation and sea-ice change. & Evidence lives in content/publications.bib and the fieldwork section. \\
Chronology and uncertainty & Making age models, proxy calibration and uncertainty legible from sample to published claim. & Evidence lives in content/publications.bib and the teaching record. \\
Open polar practice & Designing reusable datasets and reciprocal field partnerships for Southern Ocean research. & Evidence lives in the projects and service sections. \\%
}

\newcommand{\cvStrandsHeader}{%
\textbf{Title} & \textbf{Detail} & \textbf{Items} \\%
}

\newcommand{\cvStrandsInline}{%
Ocean archives (Reading marine sediments as layered records of temperature, circulation and sea-ice change.), Chronology and uncertainty (Making age models, proxy calibration and uncertainty legible from sample to published claim.), Open polar practice (Designing reusable datasets and reciprocal field partnerships for Southern Ocean research.)%
}

\newcommand{\cvStrandsCount}{3}

\newcommand{\cvStrandsPrinted}{0}

\newcommand{\cvAppointmentsNote}{%
%
}

\newcommand{\cvAppointments}{%
\ifnum1>\cvmax\else
\cventry
  {Associate Professor of Palaeoclimatology}
  {Kārearea University, School of Earth and Ice Studies}
  {Feb 2021 -- Present}
  {Ōtepoti Dunedin, Aotearoa New Zealand}
\resumeItemListStart
  \item Cross-appointed affiliate of the Tideglass Research Cooperative.
\resumeItemListEnd
\fi

\ifnum2>\cvmax\else
\cventry
  {Senior Lecturer in Marine Geology}
  {Aster Coast University, Koru School of Marine Records}
  {Jul 2016 -- Jan 2021}
  {nipaluna Hobart, Australia}
\fi%
}

\newcommand{\cvAppointmentsRows}{%
Associate Professor of Palaeoclimatology & Kārearea University & Ōtepoti Dunedin, Aotearoa New Zealand & Feb 2021 -- Present & School of Earth and Ice Studies & https://earth-ice.example.edu & Cross-appointed affiliate of the Tideglass Research Cooperative. & 2021-02-15 \\
Senior Lecturer in Marine Geology & Aster Coast University & nipaluna Hobart, Australia & Jul 2016 -- Jan 2021 & Koru School of Marine Records & 2016-07-01 \\%
}

\newcommand{\cvAppointmentsHeader}{%
\textbf{Title} & \textbf{Org} & \textbf{Place} & \textbf{Dates} & \textbf{Detail} & \textbf{Url} & \textbf{Items} & \textbf{Announced} \\%
}

\newcommand{\cvAppointmentsInline}{%
Associate Professor of Palaeoclimatology (School of Earth and Ice Studies), Senior Lecturer in Marine Geology (Koru School of Marine Records)%
}

\newcommand{\cvAppointmentsCount}{2}

\newcommand{\cvEducationNote}{%
%
}

\newcommand{\cvEducation}{%
\ifnum1>\cvmax\else
\cventry
  {Ph.D. in Marine Geology}
  {Aster Institute for Earth Records}
  {2010 -- 2014}
  {Canberra, Australia}
\resumeItemListStart
  \item Thesis: Foraminiferal Mg/Ca and Southern Ocean warming across Termination I.
\resumeItemListEnd
\fi

\ifnum2>\cvmax\else
\cventry
  {M.Sc. in Earth Sciences (First Class Honours)}
  {Kōwhai Coast University}
  {2008 -- 2009}
  {Cape Town, South Africa}
\fi

\ifnum3>\cvmax\else
\cventry
  {B.Sc. in Geology and Oceanography}
  {Tideglass Ocean Institute}
  {2005 -- 2008}
  {Chennai, India}
\fi%
}

\newcommand{\cvEducationRows}{%
Ph.D. in Marine Geology & Aster Institute for Earth Records & Canberra, Australia & 2010 -- 2014 & Thesis: Foraminiferal Mg/Ca and Southern Ocean warming across Termination I. \\
M.Sc. in Earth Sciences (First Class Honours) & Kōwhai Coast University & Cape Town, South Africa & 2008 -- 2009 \\
B.Sc. in Geology and Oceanography & Tideglass Ocean Institute & Chennai, India & 2005 -- 2008 \\%
}

\newcommand{\cvEducationHeader}{%
\textbf{Title} & \textbf{Org} & \textbf{Place} & \textbf{Dates} & \textbf{Items} \\%
}

\newcommand{\cvEducationInline}{%
Ph.D. in Marine Geology, M.Sc. in Earth Sciences (First Class Honours), B.Sc. in Geology and Oceanography%
}

\newcommand{\cvEducationCount}{3}

\newcommand{\cvTeachingNote}{%
%
}

\newcommand{\cvTeaching}{%
\ifnum1>\cvmax\else
\cventry
  {Course coordinator and lecturer}
  {Kārearea University}
  {2021 -- Present}
  {Ōtepoti Dunedin, Aotearoa New Zealand}
\cvcourses{Y Y Y Y}{\textbf{Course} & \textbf{Programme} & \textbf{Topics} & \textbf{Credit points} \\}{
GEOL 273 & B.Sc. Geology & Palaeoclimate archives; proxy calibration & 18 \\
GEOL 472 & B.Sc. (Hons) & Marine cores; age-depth modelling & 18 \\
OCEN 405 & M.Sc. Oceanography & Southern Ocean circulation; uncertainty & 20 \\}
\fi%
}

\newcommand{\cvTeachingRows}{%
Course coordinator and lecturer & Kārearea University & Ōtepoti Dunedin, Aotearoa New Zealand & 2021 -- Present & GEOL 273 B.Sc. Geology Palaeoclimate archives; proxy calibration 18, GEOL 472 B.Sc. (Hons) Marine cores; age-depth modelling 18, OCEN 405 M.Sc. Oceanography Southern Ocean circulation; uncertainty 20 \\%
}

\newcommand{\cvTeachingHeader}{%
\textbf{Title} & \textbf{Org} & \textbf{Place} & \textbf{Dates} & \textbf{Rows} \\%
}

\newcommand{\cvTeachingInline}{%
Course coordinator and lecturer%
}

\newcommand{\cvTeachingCount}{1}

\newcommand{\cvSupervisionNote}{%
\textbf{Current supervision:} projects span marine geochemistry, chronology and open polar data.%
}

\newcommand{\cvSupervision}{%
\ifnum1>\cvmax\else
\cventry
  {Doctoral candidates}
  {Primary supervisor for two; co-supervisor for two, 4}
  {}
  {}
\fi

\ifnum2>\cvmax\else
\cventry
  {M.Sc. and honours researchers}
  {Completed and current research dissertations, 9}
  {}
  {}
\fi%
}

\newcommand{\cvSupervisionRows}{%
Doctoral candidates & 4 & Primary supervisor for two; co-supervisor for two \\
M.Sc. and honours researchers & 9 & Completed and current research dissertations \\%
}

\newcommand{\cvSupervisionHeader}{%
\textbf{Title} & \textbf{Count} & \textbf{Detail} \\%
}

\newcommand{\cvSupervisionInline}{%
Doctoral candidates (Primary supervisor for two; co-supervisor for two), M.Sc. and honours researchers (Completed and current research dissertations)%
}

\newcommand{\cvSupervisionCount}{2}

\newcommand{\cvServiceNote}{%
%
}

\newcommand{\cvService}{%
\ifnum1>\cvmax\else
\cventry
  {Associate Editor}
  {Tideglass Climate Letters, Marine palaeoclimate}
  {2023 -- Present}
  {}
\fi

\ifnum2>\cvmax\else
\cventry
  {Associate Editor}
  {Example Journal of Palaeoclimate Studies, 2022--2024 \textbf{[IF: 3.4, Q1]}}
  {}
  {}
\fi

\ifnum3>\cvmax\else
\cventry
  {Scientific Steering Committee}
  {Tideglass Research Cooperative, 2022--2026}
  {}
  {}
\fi

\ifnum4>\cvmax\else
\cventry
  {Panel reviewer}
  {Kōwhai Research Trust}
  {2021 -- Present}
  {}
\fi

\ifnum5>\cvmax\else
\cventry
  {Reviewer}
  {Open Core Review}
  {}
  {}
\fi%
}

\newcommand{\cvServiceRows}{%
Associate Editor & Tideglass Climate Letters & 2023 -- Present & Marine palaeoclimate & 2023-03-10 \\
Associate Editor & Example Journal of Palaeoclimate Studies & 2022--2024 & IF: 3.4, Q1 & https://example.org/journals/palaeoclimate-studies \\
Scientific Steering Committee & Tideglass Research Cooperative & 2022--2026 \\
Panel reviewer & Kōwhai Research Trust & 2021 -- Present \\
Reviewer & Open Core Review \\%
}

\newcommand{\cvServiceHeader}{%
\textbf{Title} & \textbf{Org} & \textbf{Dates} & \textbf{Detail} & \textbf{Announced} \\%
}

\newcommand{\cvServiceInline}{%
Associate Editor (Marine palaeoclimate), Associate Editor, Scientific Steering Committee, Panel reviewer, Reviewer%
}

\newcommand{\cvServiceCount}{5}

\newcommand{\cvProjectsNote}{%
%
}

\newcommand{\cvProjects}{%
\ifnum1>\cvmax\else
\cventry
  {TIDE --- Termination Ice Dynamics and Ecosystems}
  {Kōwhai Research Trust, Principal investigator}
  {2024 -- 2027}
  {}
\resumeItemListStart
  \item Multiproxy reconstruction of Southern Ocean fronts across Termination II, with voyage support from the \href{https://voyages.example.edu}{Southern Archive Voyage Consortium}.
\resumeItemListEnd
\fi

\ifnum2>\cvmax\else
\cventry
  {Open Core Chronologies}
  {Koru Basin Research Alliance, Co-investigator}
  {2022 -- 2025}
  {}
\resumeItemListStart
  \item Reusable age-depth models and provenance for legacy sediment cores.
\resumeItemListEnd
\fi%
}

\newcommand{\cvProjectsRows}{%
TIDE --- Termination Ice Dynamics and Ecosystems & Kōwhai Research Trust & 2024 -- 2027 & Principal investigator & NZ\$960,000 & https://research.example.edu/tide & 2024-06-03 & Multiproxy reconstruction of Southern Ocean fronts across Termination II, with voyage support from the \href{https://voyages.example.edu}{Southern Archive Voyage Consortium}. \\
Open Core Chronologies & Koru Basin Research Alliance & 2022 -- 2025 & Co-investigator & A\$420,000 & Reusable age-depth models and provenance for legacy sediment cores. \\%
}

\newcommand{\cvProjectsHeader}{%
\textbf{Title} & \textbf{Org} & \textbf{Dates} & \textbf{Detail} & \textbf{Funding} & \textbf{Url} & \textbf{Announced} & \textbf{Items} \\%
}

\newcommand{\cvProjectsInline}{%
TIDE --- Termination Ice Dynamics and Ecosystems (Principal investigator), Open Core Chronologies (Co-investigator)%
}

\newcommand{\cvProjectsCount}{2}

\newcommand{\cvFieldworkNote}{%
%
}

\newcommand{\cvFieldwork}{%
\ifnum1>\cvmax\else
\cventry
  {Pukaki Rise sediment coring voyage}
  {RV Tāwhirimātea, Co-chief scientist}
  {Jan -- Feb 2025}
  {Southern Ocean}
\fi

\ifnum2>\cvmax\else
\cventry
  {East Antarctic margin core curation}
  {Koru Basin Core Archive, Visiting scientist}
  {Nov 2022}
  {Te Whanganui-a-Tara Wellington}
\fi%
}

\newcommand{\cvFieldworkRows}{%
Pukaki Rise sediment coring voyage & RV Tāwhirimātea & Southern Ocean & Jan -- Feb 2025 & Co-chief scientist & 2025-01-12 \\
East Antarctic margin core curation & Koru Basin Core Archive & Te Whanganui-a-Tara Wellington & Nov 2022 & Visiting scientist \\%
}

\newcommand{\cvFieldworkHeader}{%
\textbf{Title} & \textbf{Org} & \textbf{Place} & \textbf{Dates} & \textbf{Detail} & \textbf{Announced} \\%
}

\newcommand{\cvFieldworkInline}{%
Pukaki Rise sediment coring voyage (Co-chief scientist), East Antarctic margin core curation (Visiting scientist)%
}

\newcommand{\cvFieldworkCount}{2}

\newcommand{\cvFieldworkPrinted}{0}

\newcommand{\cvOutreachNote}{%
%
}

\newcommand{\cvOutreach}{%
\ifnum1>\cvmax\else
\cventry
  {Sediment stories: reading climate from ocean cores}
  {Kōwhai Coast Learning Collective, Community workshop for secondary-school science teachers}
  {Sep 2025}
  {Ōamaru, Aotearoa New Zealand}
\fi

\ifnum2>\cvmax\else
\cventry
  {Core archive open day}
  {Aotearoa Marine Memory Trust, Public demonstration of core scanning and age-model uncertainty}
  {Mar 2024}
  {Porirua, Aotearoa New Zealand}
\fi

\ifnum3>\cvmax\else
\cventry
  {Antarctic archives teacher residency}
  {Polar Evidence Studio, Short course on proxy evidence for secondary-school teachers}
  {Jul 2023}
  {Ōtepoti Dunedin, Aotearoa New Zealand}
\fi%
}

\newcommand{\cvOutreachRows}{%
Sediment stories: reading climate from ocean cores & Kōwhai Coast Learning Collective & Ōamaru, Aotearoa New Zealand & Sep 2025 & Community workshop for secondary-school science teachers & https://learning.example.edu/sediment-stories \\
Core archive open day & Aotearoa Marine Memory Trust & Porirua, Aotearoa New Zealand & Mar 2024 & Public demonstration of core scanning and age-model uncertainty \\
Antarctic archives teacher residency & Polar Evidence Studio & Ōtepoti Dunedin, Aotearoa New Zealand & Jul 2023 & Short course on proxy evidence for secondary-school teachers \\%
}

\newcommand{\cvOutreachHeader}{%
\textbf{Title} & \textbf{Org} & \textbf{Place} & \textbf{Dates} & \textbf{Detail} & \textbf{Url} \\%
}

\newcommand{\cvOutreachInline}{%
Sediment stories: reading climate from ocean cores (Community workshop for secondary-school science teachers), Core archive open day (Public demonstration of core scanning and age-model uncertainty), Antarctic archives teacher residency (Short course on proxy evidence for secondary-school teachers)%
}

\newcommand{\cvOutreachCount}{3}

\newcommand{\cvAwardsNote}{%
%
}

\newcommand{\cvAwards}{%
\ifnum1>\cvmax\else
\cventry
  {Early Career Research Excellence Award}
  {Kārearea Earth Forum}
  {2019}
  {}
\fi

\ifnum2>\cvmax\else
\cventry
  {Tideglass Visiting Fellowship}
  {Tideglass Research Cooperative}
  {2017}
  {Cambridge, United Kingdom}
\fi%
}

\newcommand{\cvAwardsRows}{%
Early Career Research Excellence Award & Kārearea Earth Forum & 2019 & 2019-11-20 \\
Tideglass Visiting Fellowship & Tideglass Research Cooperative & Cambridge, United Kingdom & 2017 \\%
}

\newcommand{\cvAwardsHeader}{%
\textbf{Title} & \textbf{Org} & \textbf{Dates} & \textbf{Announced} \\%
}

\newcommand{\cvAwardsInline}{%
Early Career Research Excellence Award, Tideglass Visiting Fellowship%
}

\newcommand{\cvAwardsCount}{2}

\newcommand{\cvLeadershipNote}{%
%
}

\newcommand{\cvLeadership}{%
\ifnum1>\cvmax\else
\cventry
  {Co-chair, Open Polar Data Working Group}
  {Tideglass Research Cooperative}
  {2024 -- Present}
  {}
\fi

\ifnum2>\cvmax\else
\cventry
  {Faculty representative --- \href{/blog/2026/reading-a-core/}{field note}}
  {Pacific research partnerships committee}
  {2022 -- 2024}
  {}
\fi%
}

\newcommand{\cvLeadershipRows}{%
Co-chair, Open Polar Data Working Group & Tideglass Research Cooperative & 2024 -- Present & 2024-02-01 \\
Faculty representative --- \href{/blog/2026/reading-a-core/}{field note} & Pacific research partnerships committee & 2022 -- 2024 \\%
}

\newcommand{\cvLeadershipHeader}{%
\textbf{Title} & \textbf{Org} & \textbf{Dates} & \textbf{Announced} \\%
}

\newcommand{\cvLeadershipInline}{%
Co-chair, Open Polar Data Working Group, Faculty representative --- \href{/blog/2026/reading-a-core/}{field note}%
}

\newcommand{\cvLeadershipCount}{2}

\newcommand{\cvLeadershipPrinted}{0}

\newcommand{\cvLanguagesNote}{%
%
}

\newcommand{\cvLanguages}{%
\ifnum1>\cvmax\else
\cventry
  {Tamil}
  {Native}
  {}
  {}
\fi

\ifnum2>\cvmax\else
\cventry
  {English}
  {Fluent}
  {}
  {}
\fi

\ifnum3>\cvmax\else
\cventry
  {te reo Māori}
  {Intermediate learner}
  {}
  {}
\fi%
}

\newcommand{\cvLanguagesRows}{%
Tamil & Native \\
English & Fluent \\
te reo Māori & Intermediate learner \\%
}

\newcommand{\cvLanguagesHeader}{%
\textbf{Title} & \textbf{Detail} \\%
}

\newcommand{\cvLanguagesInline}{%
Tamil (Native), English (Fluent), te reo Māori (Intermediate learner)%
}

\newcommand{\cvLanguagesCount}{3}

\newcommand{\cvPublicationsCount}{10}

\defbibfilter{Publications1}{type=article and not keyword=underreview}
\defbibfilter{Publications2}{type=inproceedings and not ( type=article and not keyword=underreview )}
\defbibfilter{Publications3}{( type=incollection or type=inbook or type=book ) and not ( type=article and not keyword=underreview ) and not ( type=inproceedings )}
\defbibfilter{Publications4}{type=dataset and not ( type=article and not keyword=underreview ) and not ( type=inproceedings ) and not ( ( type=incollection or type=inbook or type=book ) )}
\defbibfilter{Publications6}{type=article and keyword=underreview and not ( type=article and not keyword=underreview ) and not ( type=inproceedings ) and not ( ( type=incollection or type=inbook or type=book ) ) and not ( type=dataset ) and not ( type=misc )}

\newcommand{\cvPublicationsKey}{%
J=Journal, C=Conference, B=Books, D=Data, R=Under review%
}

\newcommand{\cvPublicationsSections}{%
\begin{refcontext}[labelprefix=J]
\printbibliography[heading=bibsubheading, title={Journal articles (peer-reviewed)}, filter=Publications1, resetnumbers=true]
\end{refcontext}

\begin{refcontext}[labelprefix=C]
\printbibliography[heading=bibsubheading, title={Conference papers (peer-reviewed)}, filter=Publications2, resetnumbers=true]
\end{refcontext}

\begin{refcontext}[labelprefix=B]
\printbibliography[heading=bibsubheading, title={Books and chapters}, filter=Publications3, resetnumbers=true]
\end{refcontext}

\begin{refcontext}[labelprefix=D]
\printbibliography[heading=bibsubheading, title={Data archives}, filter=Publications4, resetnumbers=true]
\end{refcontext}

\begin{refcontext}[labelprefix=R]
\printbibliography[heading=bibsubheading, title={Manuscripts under review}, filter=Publications6, resetnumbers=true]
\end{refcontext}%
}

\newcommand{\cvTalksCount}{4}

\defbibfilter{Talks1}{type=unpublished and keyword=invited}

\newcommand{\cvTalksKey}{%
I=Invited%
}

\newcommand{\cvTalksSections}{%
\begin{refcontext}[labelprefix=I]
\printbibliography[heading=bibsubheading, title={Invited and keynote talks}, filter=Talks1, resetnumbers=true]
\end{refcontext}%
}

\newcommand{\cvAutoSections}{%
\cvautopart{Teaching Statement}{TeachingStatement}%
\cvautopart{Strands}{Strands}%
\cvautopart{Appointments}{Appointments}%
\cvautopart{Education}{Education}%
\cvautopart{Teaching}{Teaching}%
\cvautopart{Supervision}{Supervision}%
\cvautopart{Service}{Service}%
\cvautopart{Projects}{Projects}%
\cvautopart{Fieldwork}{Fieldwork}%
\cvautopart{Public engagement}{Outreach}%
\cvautopart{Awards}{Awards}%
\cvautopart{Leadership}{Leadership}%
\cvautopart{Languages}{Languages}%
}


,
% so a change to the printed design is made once and reaches all of them.
%
% To change what a CV says, edit content/cv.yaml and regenerate. To change how
% it looks, edit this file. To change WHICH of it one document prints, edit that
% document's body - cv/cv.tex, cv/short.tex or cv/teaching.tex.
%
% WHICH sections print is a record fact, not a layout one. The `\cvpart` lines in
% a document body places the sections whose heading or setting that document
% curates; `\cvAutoSections` then prints every remaining section of
% content/cv.yaml, in the order the record writes them, under the heading its key
% spells out. Adding a section to the record therefore needs no edit here. Give
% a section `printed: false` in content/cv.yaml to keep it off the PDF - that
% holds for the curated ones below too - and `heading:` to name an uncurated one
% something other than its key.
%
% Build:  latexmk -xelatex -cd cv/cv.tex      (xelatex is required: academicons)
%         the same for cv/short.tex and cv/teaching.tex
% -----------------------------------------------------------------------------



\documentclass[a4paper,11pt]{article}

% Package imports
\usepackage{latexsym}
\usepackage{xcolor}
\usepackage{float}
\usepackage{ragged2e}
\usepackage[empty]{fullpage}
\usepackage{wrapfig}
\usepackage{lipsum}
\usepackage{booktabs}   % high-quality horizontal rules (no vertical rules)
\usepackage{array}      % better column types
\usepackage{xparse}
\usepackage{tabularx}
\usepackage{titlesec}
\usepackage{geometry}
\usepackage{marvosym}
\usepackage{verbatim}
\usepackage{enumitem}
\usepackage{fancyhdr}
\usepackage{multicol}
\usepackage{graphicx}
\usepackage{fontspec}
\setmainfont{TeX Gyre Pagella}
\usepackage{fontawesome5}
\usepackage{academicons} % <-- NEW (Scholar, ORCID, etc.)
\usepackage{eurosym} % <-- NEW (€)

% Column helpers (tight but readable)
\newcolumntype{P}[1]{>{\raggedright\arraybackslash}p{#1}}
\newcolumntype{Y}{>{\raggedright\arraybackslash}X}

% A reusable CV table environment (no floats, correct spacing)
\NewDocumentEnvironment{cvtable}{m +b}{%
  \vspace{0.25em}
  \small
  \setlength{\tabcolsep}{5pt}
  \renewcommand{\arraystretch}{1.15}
  \begin{tabularx}{\linewidth}{@{}#1@{}}
  #2
  \end{tabularx}
  \vspace{0.25em}
}{}

% Color definitions
\definecolor{darkblue}{RGB}{0,0,139}

% Page layout
\setlength{\multicolsep}{0pt}
\pagestyle{fancy}
\fancyhf{} % clear all header and footer fields
\fancyfoot{}
\renewcommand{\headrulewidth}{0pt}
\renewcommand{\footrulewidth}{0pt}
\geometry{left=1.4cm, top=0.8cm, right=1.2cm, bottom=1cm}
\setlength{\footskip}{5pt} % Addressing fancyhdr warning

% Hyperlink setup (moved after fancyhdr to address warning)
\usepackage[hidelinks]{hyperref}
\hypersetup{
    colorlinks=true,
    linkcolor=darkblue,
    filecolor=darkblue,
    urlcolor=darkblue,
}

% Custom box settings
\usepackage[most]{tcolorbox}
\tcbset{
    frame code={},
    center title,
    left=0pt,
    right=0pt,
    top=0pt,
    bottom=0pt,
    colback=gray!20,
    colframe=white,
    width=\dimexpr\textwidth\relax,
    enlarge left by=-2mm,
    boxsep=4pt,
    arc=0pt,outer arc=0pt,
}

% URL style
\urlstyle{same}

% Text alignment
\raggedright
\setlength{\tabcolsep}{0in}

% Section formatting
\titleformat{\section}{
  \vspace{-4pt}\scshape\raggedright\large
}{}{0em}{}[\color{black}\titlerule \vspace{-7pt}]

% Custom commands
\newcommand{\resumeItem}[2]{
  \item{
    \textbf{#1}{\hspace{0.5mm}#2 \vspace{-0.5mm}}
  }
}

\newcommand{\resumePOR}[3]{
\vspace{0.5mm}\item
    \begin{tabular*}{0.97\textwidth}[t]{l@{\extracolsep{\fill}}r}
        \textbf{#1}\hspace{0.3mm}#2 & \textit{\small{#3}}
    \end{tabular*}
    \vspace{-2mm}
}

% One entry, every section: what it was, where it was, when, and the place.
% `content/cv.yaml` has exactly one entry shape and this is how it is set.
\newcommand{\cventry}[4]{
\vspace{0.5mm}\item
    \begin{tabularx}{0.98\textwidth}[t]{@{}Y >{\raggedleft\arraybackslash}p{4.8cm}@{}}
        \textbf{#1} & \textit{\footnotesize{#3}} \\
        \textit{\footnotesize{#2}} &  \footnotesize{#4}\\
    \end{tabularx}
    \vspace{-2.4mm}
}

% An empty generated macro, to test the optional prose fields against.
% NOT `\ifx\cvFocus\empty`: \newcommand makes its macros \long and \empty is
% not, so \ifx compares false for two macros that both hold nothing, and the
% heading prints over an empty field. This one is made the same way they are.
\newcommand{\cvempty}{}

% Separates the paragraphs of a section's `note:`. Spacing is layout, so the
% amount of it lives here and the generator only says where the break goes.
\newcommand{\cvnotesep}{\par\vspace{1.2mm}}

% Record that this file lays out one section key itself, so \cvautopart below
% does not print it a second time from the record's own section sequence.
% \gdef, not \def: \cvpartflush sets its entries inside a group.
\newcommand{\cvlaidout}[1]{\expandafter\gdef\csname cvlaidout@#1\endcsname{}}

% Make the five generated macros this layout may mention safe even when
% content/cv.yaml has no such section, and default the section to printing.
% (`Header` is the sixth generated macro, but this layout supplies its own table
% headings.) `\cvpart` does this for you; a hand-written section block must call it
% itself.
\newcommand{\cvdeclare}[1]{%
  \cvlaidout{#1}%
  % Printing is the default: only `printed: false` in content/cv.yaml generates a
  % \cv<Key>Printed of 0, and this \providecommand leaves that definition alone.
  \expandafter\providecommand\csname cv#1Printed\endcsname{1}%
  \expandafter\providecommand\csname cv#1Count\endcsname{0}%
  \expandafter\providecommand\csname cv#1\endcsname{}%
  \expandafter\providecommand\csname cv#1Note\endcsname{}%
  \expandafter\providecommand\csname cv#1Rows\endcsname{}%
  \expandafter\providecommand\csname cv#1Inline\endcsname{}%
}

% The same, for a bibliography section: `publications:`/`talks:` in cv.yaml
% generate \cv<Name>Key, \cv<Name>Sections and \cv<Name>Count (the number of
% entries in the matching .bib) instead of the five above. An adopter who has
% declared neither - which is everyone on day one - would otherwise meet a raw
% "Undefined control sequence" that stops xelatex at its interactive prompt.
% Empty is the documented behaviour: the heading prints over an empty
% bibliography and README.md tells them to delete the block.
\newcommand{\cvdeclarebib}[1]{%
  \expandafter\providecommand\csname cv#1Count\endcsname{0}%
  \expandafter\providecommand\csname cv#1Key\endcsname{}%
  \expandafter\providecommand\csname cv#1Sections\endcsname{}%
}

% How many entries of a section a document prints. The generated \cv<Name>
% guards each of its entries with \ifnum<n>>\cvmax, so this is the cutoff.
%
% This is the layout's call, and it is a different question from the record's
% `printed:`. WHICH sections print is a record fact; HOW MANY entries of one a
% document prints is a page budget belonging to that document - the same section
% is worth three entries in a short CV and all of them in the full one - so the
% number lives here and content/cv.yaml never states it. A count high enough
% that no record reaches it is the default, which keeps the full CV whole.
%
% Curation is neither: "my five best papers" is a judgement about the work, so
% it is marked on the fact itself - see the `selected` keyword in
% content/publications.bib and how cv/short.tex filters on it.
\newcommand{\cvmax}{99999}

% Print one section of content/cv.yaml. #2 is the heading, #3 the section's
% macro name (`Awards` for the `awards:` key), and the optional #1 is how many
% of its entries to print. A section the file does not have, or gives
% `printed: false`, prints nothing at all - no heading, no gap - so this line
% can stay. Reorder the CV by reordering these lines in the document body.
%
% The whole thing is one group so that \cvmax is restored afterwards: a
% \cvpart[3] must not silently truncate the section written after it.
\newcommand{\cvpart}[3][99999]{{%
  \renewcommand{\cvmax}{#1}%
  \cvdeclare{#3}%
  \ifnum\csname cv#3Count\endcsname>0
  \ifnum\csname cv#3Printed\endcsname>0
  \section{\textbf{#2}}
  \vspace{-0.4mm}
  \resumeSubHeadingListStart
  \csname cv#3\endcsname
  \resumeSubHeadingListEnd
  \vspace{-6mm}
  \fi
  \fi
}}

% One section of the record's own section sequence: `\cvpart`, unless this file
% has already laid that key out by hand above. Every line of \cvAutoSections is
% one of these, so an adopter's new section prints with no edit here, and the
% ones curated below keep their curated heading, order and setting.
\newcommand{\cvautopart}[2]{%
  \expandafter\ifx\csname cvlaidout@#2\endcsname\relax\cvpart{#1}{#2}\fi%
}

% `\cvpart`, with the second line of every entry set the other way. Which
% sections use which is a layout choice and lives here; cv.yaml says nothing
% about it.
\newcommand{\cvpartflush}[3][99999]{{%
  \renewcommand{\cventry}[4]{%
    \vspace{0.5mm}\item
    \begin{tabularx}{0.98\textwidth}[t]{@{}Y >{\raggedleft\arraybackslash}p{4.8cm}@{}}
        \textbf{##1} & \textit{\footnotesize{##3}} \\
        \footnotesize{\textit{##2}} & \footnotesize{##4}
    \end{tabularx}
    \vspace{-2.4mm}
  }%
  \cvpart[#1]{#2}{#3}}}

% A table hanging under an entry: `rows:` in cv.yaml. #1 is a generated
% equal-width column specification, #2 the generated headings and #3 the rows.
% Any number of consistently shaped columns therefore works without editing
% this layout. The generator refuses rows whose key order disagrees.
\newcommand{\cvcourses}[3]{%
\par
\begin{cvtable}{#1}
\toprule
#2
\midrule
#3
\bottomrule
\end{cvtable}
}

\newcommand{\resumeSubItem}[2]{\resumeItem{#1}{#2}\vspace{-4pt}}

\renewcommand{\labelitemi}{$\vcenter{\hbox{\tiny$\bullet$}}$}
\renewcommand{\labelitemii}{$\vcenter{\hbox{\tiny$\circ$}}$}

\newcommand{\resumeSubHeadingListStart}{\begin{itemize}[leftmargin=*,labelsep=1mm]}
\newcommand{\resumeHeadingSkillStart}{\begin{itemize}[leftmargin=*,itemsep=1.7mm, rightmargin=2ex]}
\newcommand{\resumeItemListStart}{\begin{itemize}[leftmargin=*,labelsep=1mm,itemsep=0.5mm]}

\newcommand{\resumeSubHeadingListEnd}{\end{itemize}\vspace{2mm}}
\newcommand{\resumeHeadingSkillEnd}{\end{itemize}\vspace{-2mm}}
\newcommand{\resumeItemListEnd}{\end{itemize}\vspace{-2mm}}
\newcommand{\cvsection}[1]{%
\vspace{2mm}
\begin{tcolorbox}
    \textbf{\large #1}
\end{tcolorbox}
    \vspace{-4mm}
}

\newcolumntype{L}{>{\raggedright\arraybackslash}X}%
\newcolumntype{R}{>{\raggedleft\arraybackslash}X}%
\newcolumntype{C}{>{\centering\arraybackslash}X}%

% Commands for icon sizing and positioning
\newcommand{\socialicon}[1]{\raisebox{-0.05em}{\resizebox{!}{1em}{#1}}}
\newcommand{\ieeeicon}[1]{\raisebox{-0.3em}{\resizebox{!}{1.3em}{#1}}}

% Profile icon registry. Known compact account IDs select their service icon;
% arbitrary `{label, url}` links use \cviconlink.
\newcommand{\cviconscholar}{\socialicon{\aiGoogleScholar}}
\newcommand{\cviconorcid}{\socialicon{\aiOrcid}}
\newcommand{\cviconlinkedin}{\socialicon{\faLinkedin}}
\newcommand{\cvicongithub}{\socialicon{\faGithub}}
\newcommand{\cviconlink}{\socialicon{\faLink}}

% Font options
\newcommand{\headerfonti}{\fontfamily{phv}\selectfont} % Helvetica-like (similar to Arial/Calibri)
\newcommand{\headerfontii}{\fontfamily{ptm}\selectfont} % Times-like (similar to Times New Roman)
\newcommand{\headerfontiii}{\rmfamily} % TeX Gyre Pagella, set above with Unicode support
\newcommand{\headerfontiv}{\fontfamily{pbk}\selectfont} % Bookman (readable serif)
\newcommand{\headerfontv}{\fontfamily{pag}\selectfont} % Avant Garde-like (similar to Trebuchet MS)
\newcommand{\headerfontvi}{\fontfamily{cmss}\selectfont} % Computer Modern Sans Serif
\newcommand{\headerfontvii}{\fontfamily{qhv}\selectfont} % Quasi-Helvetica (another Arial/Calibri alternative)
\newcommand{\headerfontviii}{\fontfamily{qpl}\selectfont} % Quasi-Palatino (another elegant serif option)
\newcommand{\headerfontix}{\fontfamily{qtm}\selectfont} % Quasi-Times (another Times New Roman alternative)
\newcommand{\headerfontx}{\fontfamily{bch}\selectfont} % Charter (clean serif font)

% --------------------------------
% BIBLATEX (publications + talks)
% --------------------------------
\usepackage[
  backend=biber,
  style=numeric,
  sorting=ydnt,
  defernumbers=true,
  maxbibnames=99,
  giveninits=true,
  doi=true,
  url=true,
  isbn=false
]{biblatex}

% Two bib databases (kept separate via refsection below).
% Publications are the repository-wide canonical list, shared with the website.
\addbibresource{../content/publications.bib}
\addbibresource{../content/talks.bib}

% The shared bibliography may carry rendering fields used only by the website. They are not part of the CV, and abstracts can contain raw % and &
% that would break the LaTeX pass, so drop them during biber processing instead
% of altering the .bib the website reads.
\DeclareSourcemap{
  \maps[datatype=bibtex]{
    \map{
      \step[fieldset=abstract, null]
      \step[fieldset=selected, null]
      \step[fieldset=preview, null]
      \step[fieldset=abbr, null]
      \step[fieldset=bibtex_show, null]
      \step[fieldset=html, null]
      \step[fieldset=pdf, null]
      \step[fieldset=code, null]
      \step[fieldset=altmetric, null]
      \step[fieldset=dimensions, null]
    }
    % DBLP's @misc artifact records carry an UNFILLED placeholder,
    % `note = {Accessed on YYYY-MM-DD.}`, which biblatex would print verbatim as
    % the venue of a released software package. The website already skips it
    % (`VENUE_FIELDS` in web/src/lib/record.ts prefers `publisher`); this is the
    % same fix on the printed side, so declaring a `types: [misc]` section in
    % content/cv.yaml renders correctly instead of printing the placeholder.
    \map{
      \step[fieldsource=note, match=\regexp{\AAccessed\son\sYYYY-MM-DD\.?\Z}, final]
      \step[fieldset=note, null]
    }
  }
}

% Compact bib
\setlength\bibitemsep{0.25\baselineskip}
\renewcommand*{\bibfont}{\footnotesize}

% Bib headings matching your tidy style
\defbibheading{bibsubheading}[\refname]{%
  \vspace{0.8mm}\textbf{#1}\par\vspace{-0.3mm}
}

% No filters here. Which entry type belongs under which heading is a curated
% opinion, so it is declared in `content/cv.yaml` under `publications:` and
% `talks:` and read by the website too. build-cv-data.mjs turns each declared
% section into its \defbibfilter and its \printbibliography, and gives this file
% \cvPublicationsKey / \cvPublicationsSections and \cvTalksKey /
% \cvTalksSections, plus \cvPublicationsCount / \cvTalksCount - how many entries
% the .bib holds, which guards the \refsection below.
% Adding, reordering or renaming a section is a cv.yaml edit.

% --------------------------------
% CONTENT (generated from content/cv.yaml)
% --------------------------------
% =============================================================================
% GENERATED FILE - DO NOT EDIT.
%
% Produced by `node scripts/build-cv-data.mjs` from `content/cv.yaml`.
% Edit cv.yaml and regenerate; hand edits here are overwritten and CI rejects
% them (the workflow fails if this file is stale relative to cv.yaml).
%
% This file holds CONTENT ONLY. Layout, spacing and styling live in cv/preamble.tex.
% =============================================================================

\newcommand{\cvName}{%
Sahana Aster KŌWHAI, Ph.D.%
}

\newcommand{\cvContactLine}{%
\href{mailto:sahana.kowhai@example.edu}{sahana.kowhai@example.edu}
\;|\;
\href{https://profiles.example.edu/sahana-kowhai}{Staff profile}%
}

\newcommand{\cvProfilesLine}{%
\cviconorcid\,\href{https://orcid.org/0000-0000-0000-0000}{0000-0000-0000-0000}
\;|\;
\cviconscholar\,\href{https://scholar.google.com/citations?user=sahana-kowhai-example}{sahana-kowhai-example}
\;|\;
\cvicongithub\,\href{https://github.com/sahana-kowhai}{sahana-kowhai}
\;|\;
\cviconlinkedin\,\href{https://www.linkedin.com/in/sahana-kowhai}{sahana-kowhai}
\;|\;
\cviconlink\,\href{https://bsky.app/profile/sahana-kowhai.example}{Bluesky}%
}

\newcommand{\cvAffiliation}{%
Associate Professor of Palaeoclimatology, School of Earth and Ice Studies \\
Kārearea University, \\
Tideglass Research Cooperative, Ōtepoti Dunedin, Aotearoa New Zealand%
}

\newcommand{\cvShortBio}{%
Sahana Aster Kōwhai is a \textbf{fictional} palaeoclimatologist whose \emph{research} reconstructs Southern Ocean change from marine sediment archives.%
}

\newcommand{\cvFocus}{%
Southern Ocean palaeoceanography across glacial terminations, with an emphasis on multiproxy reconstructions, reproducible chronology and equitable Antarctic field partnerships.%
}

\newcommand{\cvFooter}{%
\textbf{Example record.} This curriculum vitae describes a fictional scholar and is supplied only to demonstrate the template.%
}

\newcommand{\cvTeachingStatementNote}{%
I teach palaeoclimate reasoning as \textbf{evidence work}: students leave able to defend an age model, not only to run one.%
}

\newcommand{\cvTeachingStatement}{%
%
}

\newcommand{\cvTeachingStatementRows}{%
%
}

\newcommand{\cvTeachingStatementHeader}{%
%
}

\newcommand{\cvTeachingStatementInline}{%
%
}

\newcommand{\cvTeachingStatementCount}{0}

\newcommand{\cvStrandsNote}{%
%
}

\newcommand{\cvStrands}{%
\ifnum1>\cvmax\else
\cventry
  {Ocean archives}
  {Reading marine sediments as layered records of temperature, circulation and sea-ice change.}
  {}
  {}
\resumeItemListStart
  \item Evidence lives in content/publications.bib and the fieldwork section.
\resumeItemListEnd
\fi

\ifnum2>\cvmax\else
\cventry
  {Chronology and uncertainty}
  {Making age models, proxy calibration and uncertainty legible from sample to published claim.}
  {}
  {}
\resumeItemListStart
  \item Evidence lives in content/publications.bib and the teaching record.
\resumeItemListEnd
\fi

\ifnum3>\cvmax\else
\cventry
  {Open polar practice}
  {Designing reusable datasets and reciprocal field partnerships for Southern Ocean research.}
  {}
  {}
\resumeItemListStart
  \item Evidence lives in the projects and service sections.
\resumeItemListEnd
\fi%
}

\newcommand{\cvStrandsRows}{%
Ocean archives & Reading marine sediments as layered records of temperature, circulation and sea-ice change. & Evidence lives in content/publications.bib and the fieldwork section. \\
Chronology and uncertainty & Making age models, proxy calibration and uncertainty legible from sample to published claim. & Evidence lives in content/publications.bib and the teaching record. \\
Open polar practice & Designing reusable datasets and reciprocal field partnerships for Southern Ocean research. & Evidence lives in the projects and service sections. \\%
}

\newcommand{\cvStrandsHeader}{%
\textbf{Title} & \textbf{Detail} & \textbf{Items} \\%
}

\newcommand{\cvStrandsInline}{%
Ocean archives (Reading marine sediments as layered records of temperature, circulation and sea-ice change.), Chronology and uncertainty (Making age models, proxy calibration and uncertainty legible from sample to published claim.), Open polar practice (Designing reusable datasets and reciprocal field partnerships for Southern Ocean research.)%
}

\newcommand{\cvStrandsCount}{3}

\newcommand{\cvStrandsPrinted}{0}

\newcommand{\cvAppointmentsNote}{%
%
}

\newcommand{\cvAppointments}{%
\ifnum1>\cvmax\else
\cventry
  {Associate Professor of Palaeoclimatology}
  {Kārearea University, School of Earth and Ice Studies}
  {Feb 2021 -- Present}
  {Ōtepoti Dunedin, Aotearoa New Zealand}
\resumeItemListStart
  \item Cross-appointed affiliate of the Tideglass Research Cooperative.
\resumeItemListEnd
\fi

\ifnum2>\cvmax\else
\cventry
  {Senior Lecturer in Marine Geology}
  {Aster Coast University, Koru School of Marine Records}
  {Jul 2016 -- Jan 2021}
  {nipaluna Hobart, Australia}
\fi%
}

\newcommand{\cvAppointmentsRows}{%
Associate Professor of Palaeoclimatology & Kārearea University & Ōtepoti Dunedin, Aotearoa New Zealand & Feb 2021 -- Present & School of Earth and Ice Studies & https://earth-ice.example.edu & Cross-appointed affiliate of the Tideglass Research Cooperative. & 2021-02-15 \\
Senior Lecturer in Marine Geology & Aster Coast University & nipaluna Hobart, Australia & Jul 2016 -- Jan 2021 & Koru School of Marine Records & 2016-07-01 \\%
}

\newcommand{\cvAppointmentsHeader}{%
\textbf{Title} & \textbf{Org} & \textbf{Place} & \textbf{Dates} & \textbf{Detail} & \textbf{Url} & \textbf{Items} & \textbf{Announced} \\%
}

\newcommand{\cvAppointmentsInline}{%
Associate Professor of Palaeoclimatology (School of Earth and Ice Studies), Senior Lecturer in Marine Geology (Koru School of Marine Records)%
}

\newcommand{\cvAppointmentsCount}{2}

\newcommand{\cvEducationNote}{%
%
}

\newcommand{\cvEducation}{%
\ifnum1>\cvmax\else
\cventry
  {Ph.D. in Marine Geology}
  {Aster Institute for Earth Records}
  {2010 -- 2014}
  {Canberra, Australia}
\resumeItemListStart
  \item Thesis: Foraminiferal Mg/Ca and Southern Ocean warming across Termination I.
\resumeItemListEnd
\fi

\ifnum2>\cvmax\else
\cventry
  {M.Sc. in Earth Sciences (First Class Honours)}
  {Kōwhai Coast University}
  {2008 -- 2009}
  {Cape Town, South Africa}
\fi

\ifnum3>\cvmax\else
\cventry
  {B.Sc. in Geology and Oceanography}
  {Tideglass Ocean Institute}
  {2005 -- 2008}
  {Chennai, India}
\fi%
}

\newcommand{\cvEducationRows}{%
Ph.D. in Marine Geology & Aster Institute for Earth Records & Canberra, Australia & 2010 -- 2014 & Thesis: Foraminiferal Mg/Ca and Southern Ocean warming across Termination I. \\
M.Sc. in Earth Sciences (First Class Honours) & Kōwhai Coast University & Cape Town, South Africa & 2008 -- 2009 \\
B.Sc. in Geology and Oceanography & Tideglass Ocean Institute & Chennai, India & 2005 -- 2008 \\%
}

\newcommand{\cvEducationHeader}{%
\textbf{Title} & \textbf{Org} & \textbf{Place} & \textbf{Dates} & \textbf{Items} \\%
}

\newcommand{\cvEducationInline}{%
Ph.D. in Marine Geology, M.Sc. in Earth Sciences (First Class Honours), B.Sc. in Geology and Oceanography%
}

\newcommand{\cvEducationCount}{3}

\newcommand{\cvTeachingNote}{%
%
}

\newcommand{\cvTeaching}{%
\ifnum1>\cvmax\else
\cventry
  {Course coordinator and lecturer}
  {Kārearea University}
  {2021 -- Present}
  {Ōtepoti Dunedin, Aotearoa New Zealand}
\cvcourses{Y Y Y Y}{\textbf{Course} & \textbf{Programme} & \textbf{Topics} & \textbf{Credit points} \\}{
GEOL 273 & B.Sc. Geology & Palaeoclimate archives; proxy calibration & 18 \\
GEOL 472 & B.Sc. (Hons) & Marine cores; age-depth modelling & 18 \\
OCEN 405 & M.Sc. Oceanography & Southern Ocean circulation; uncertainty & 20 \\}
\fi%
}

\newcommand{\cvTeachingRows}{%
Course coordinator and lecturer & Kārearea University & Ōtepoti Dunedin, Aotearoa New Zealand & 2021 -- Present & GEOL 273 B.Sc. Geology Palaeoclimate archives; proxy calibration 18, GEOL 472 B.Sc. (Hons) Marine cores; age-depth modelling 18, OCEN 405 M.Sc. Oceanography Southern Ocean circulation; uncertainty 20 \\%
}

\newcommand{\cvTeachingHeader}{%
\textbf{Title} & \textbf{Org} & \textbf{Place} & \textbf{Dates} & \textbf{Rows} \\%
}

\newcommand{\cvTeachingInline}{%
Course coordinator and lecturer%
}

\newcommand{\cvTeachingCount}{1}

\newcommand{\cvSupervisionNote}{%
\textbf{Current supervision:} projects span marine geochemistry, chronology and open polar data.%
}

\newcommand{\cvSupervision}{%
\ifnum1>\cvmax\else
\cventry
  {Doctoral candidates}
  {Primary supervisor for two; co-supervisor for two, 4}
  {}
  {}
\fi

\ifnum2>\cvmax\else
\cventry
  {M.Sc. and honours researchers}
  {Completed and current research dissertations, 9}
  {}
  {}
\fi%
}

\newcommand{\cvSupervisionRows}{%
Doctoral candidates & 4 & Primary supervisor for two; co-supervisor for two \\
M.Sc. and honours researchers & 9 & Completed and current research dissertations \\%
}

\newcommand{\cvSupervisionHeader}{%
\textbf{Title} & \textbf{Count} & \textbf{Detail} \\%
}

\newcommand{\cvSupervisionInline}{%
Doctoral candidates (Primary supervisor for two; co-supervisor for two), M.Sc. and honours researchers (Completed and current research dissertations)%
}

\newcommand{\cvSupervisionCount}{2}

\newcommand{\cvServiceNote}{%
%
}

\newcommand{\cvService}{%
\ifnum1>\cvmax\else
\cventry
  {Associate Editor}
  {Tideglass Climate Letters, Marine palaeoclimate}
  {2023 -- Present}
  {}
\fi

\ifnum2>\cvmax\else
\cventry
  {Associate Editor}
  {Example Journal of Palaeoclimate Studies, 2022--2024 \textbf{[IF: 3.4, Q1]}}
  {}
  {}
\fi

\ifnum3>\cvmax\else
\cventry
  {Scientific Steering Committee}
  {Tideglass Research Cooperative, 2022--2026}
  {}
  {}
\fi

\ifnum4>\cvmax\else
\cventry
  {Panel reviewer}
  {Kōwhai Research Trust}
  {2021 -- Present}
  {}
\fi

\ifnum5>\cvmax\else
\cventry
  {Reviewer}
  {Open Core Review}
  {}
  {}
\fi%
}

\newcommand{\cvServiceRows}{%
Associate Editor & Tideglass Climate Letters & 2023 -- Present & Marine palaeoclimate & 2023-03-10 \\
Associate Editor & Example Journal of Palaeoclimate Studies & 2022--2024 & IF: 3.4, Q1 & https://example.org/journals/palaeoclimate-studies \\
Scientific Steering Committee & Tideglass Research Cooperative & 2022--2026 \\
Panel reviewer & Kōwhai Research Trust & 2021 -- Present \\
Reviewer & Open Core Review \\%
}

\newcommand{\cvServiceHeader}{%
\textbf{Title} & \textbf{Org} & \textbf{Dates} & \textbf{Detail} & \textbf{Announced} \\%
}

\newcommand{\cvServiceInline}{%
Associate Editor (Marine palaeoclimate), Associate Editor, Scientific Steering Committee, Panel reviewer, Reviewer%
}

\newcommand{\cvServiceCount}{5}

\newcommand{\cvProjectsNote}{%
%
}

\newcommand{\cvProjects}{%
\ifnum1>\cvmax\else
\cventry
  {TIDE --- Termination Ice Dynamics and Ecosystems}
  {Kōwhai Research Trust, Principal investigator}
  {2024 -- 2027}
  {}
\resumeItemListStart
  \item Multiproxy reconstruction of Southern Ocean fronts across Termination II, with voyage support from the \href{https://voyages.example.edu}{Southern Archive Voyage Consortium}.
\resumeItemListEnd
\fi

\ifnum2>\cvmax\else
\cventry
  {Open Core Chronologies}
  {Koru Basin Research Alliance, Co-investigator}
  {2022 -- 2025}
  {}
\resumeItemListStart
  \item Reusable age-depth models and provenance for legacy sediment cores.
\resumeItemListEnd
\fi%
}

\newcommand{\cvProjectsRows}{%
TIDE --- Termination Ice Dynamics and Ecosystems & Kōwhai Research Trust & 2024 -- 2027 & Principal investigator & NZ\$960,000 & https://research.example.edu/tide & 2024-06-03 & Multiproxy reconstruction of Southern Ocean fronts across Termination II, with voyage support from the \href{https://voyages.example.edu}{Southern Archive Voyage Consortium}. \\
Open Core Chronologies & Koru Basin Research Alliance & 2022 -- 2025 & Co-investigator & A\$420,000 & Reusable age-depth models and provenance for legacy sediment cores. \\%
}

\newcommand{\cvProjectsHeader}{%
\textbf{Title} & \textbf{Org} & \textbf{Dates} & \textbf{Detail} & \textbf{Funding} & \textbf{Url} & \textbf{Announced} & \textbf{Items} \\%
}

\newcommand{\cvProjectsInline}{%
TIDE --- Termination Ice Dynamics and Ecosystems (Principal investigator), Open Core Chronologies (Co-investigator)%
}

\newcommand{\cvProjectsCount}{2}

\newcommand{\cvFieldworkNote}{%
%
}

\newcommand{\cvFieldwork}{%
\ifnum1>\cvmax\else
\cventry
  {Pukaki Rise sediment coring voyage}
  {RV Tāwhirimātea, Co-chief scientist}
  {Jan -- Feb 2025}
  {Southern Ocean}
\fi

\ifnum2>\cvmax\else
\cventry
  {East Antarctic margin core curation}
  {Koru Basin Core Archive, Visiting scientist}
  {Nov 2022}
  {Te Whanganui-a-Tara Wellington}
\fi%
}

\newcommand{\cvFieldworkRows}{%
Pukaki Rise sediment coring voyage & RV Tāwhirimātea & Southern Ocean & Jan -- Feb 2025 & Co-chief scientist & 2025-01-12 \\
East Antarctic margin core curation & Koru Basin Core Archive & Te Whanganui-a-Tara Wellington & Nov 2022 & Visiting scientist \\%
}

\newcommand{\cvFieldworkHeader}{%
\textbf{Title} & \textbf{Org} & \textbf{Place} & \textbf{Dates} & \textbf{Detail} & \textbf{Announced} \\%
}

\newcommand{\cvFieldworkInline}{%
Pukaki Rise sediment coring voyage (Co-chief scientist), East Antarctic margin core curation (Visiting scientist)%
}

\newcommand{\cvFieldworkCount}{2}

\newcommand{\cvFieldworkPrinted}{0}

\newcommand{\cvOutreachNote}{%
%
}

\newcommand{\cvOutreach}{%
\ifnum1>\cvmax\else
\cventry
  {Sediment stories: reading climate from ocean cores}
  {Kōwhai Coast Learning Collective, Community workshop for secondary-school science teachers}
  {Sep 2025}
  {Ōamaru, Aotearoa New Zealand}
\fi

\ifnum2>\cvmax\else
\cventry
  {Core archive open day}
  {Aotearoa Marine Memory Trust, Public demonstration of core scanning and age-model uncertainty}
  {Mar 2024}
  {Porirua, Aotearoa New Zealand}
\fi

\ifnum3>\cvmax\else
\cventry
  {Antarctic archives teacher residency}
  {Polar Evidence Studio, Short course on proxy evidence for secondary-school teachers}
  {Jul 2023}
  {Ōtepoti Dunedin, Aotearoa New Zealand}
\fi%
}

\newcommand{\cvOutreachRows}{%
Sediment stories: reading climate from ocean cores & Kōwhai Coast Learning Collective & Ōamaru, Aotearoa New Zealand & Sep 2025 & Community workshop for secondary-school science teachers & https://learning.example.edu/sediment-stories \\
Core archive open day & Aotearoa Marine Memory Trust & Porirua, Aotearoa New Zealand & Mar 2024 & Public demonstration of core scanning and age-model uncertainty \\
Antarctic archives teacher residency & Polar Evidence Studio & Ōtepoti Dunedin, Aotearoa New Zealand & Jul 2023 & Short course on proxy evidence for secondary-school teachers \\%
}

\newcommand{\cvOutreachHeader}{%
\textbf{Title} & \textbf{Org} & \textbf{Place} & \textbf{Dates} & \textbf{Detail} & \textbf{Url} \\%
}

\newcommand{\cvOutreachInline}{%
Sediment stories: reading climate from ocean cores (Community workshop for secondary-school science teachers), Core archive open day (Public demonstration of core scanning and age-model uncertainty), Antarctic archives teacher residency (Short course on proxy evidence for secondary-school teachers)%
}

\newcommand{\cvOutreachCount}{3}

\newcommand{\cvAwardsNote}{%
%
}

\newcommand{\cvAwards}{%
\ifnum1>\cvmax\else
\cventry
  {Early Career Research Excellence Award}
  {Kārearea Earth Forum}
  {2019}
  {}
\fi

\ifnum2>\cvmax\else
\cventry
  {Tideglass Visiting Fellowship}
  {Tideglass Research Cooperative}
  {2017}
  {Cambridge, United Kingdom}
\fi%
}

\newcommand{\cvAwardsRows}{%
Early Career Research Excellence Award & Kārearea Earth Forum & 2019 & 2019-11-20 \\
Tideglass Visiting Fellowship & Tideglass Research Cooperative & Cambridge, United Kingdom & 2017 \\%
}

\newcommand{\cvAwardsHeader}{%
\textbf{Title} & \textbf{Org} & \textbf{Dates} & \textbf{Announced} \\%
}

\newcommand{\cvAwardsInline}{%
Early Career Research Excellence Award, Tideglass Visiting Fellowship%
}

\newcommand{\cvAwardsCount}{2}

\newcommand{\cvLeadershipNote}{%
%
}

\newcommand{\cvLeadership}{%
\ifnum1>\cvmax\else
\cventry
  {Co-chair, Open Polar Data Working Group}
  {Tideglass Research Cooperative}
  {2024 -- Present}
  {}
\fi

\ifnum2>\cvmax\else
\cventry
  {Faculty representative --- \href{/blog/2026/reading-a-core/}{field note}}
  {Pacific research partnerships committee}
  {2022 -- 2024}
  {}
\fi%
}

\newcommand{\cvLeadershipRows}{%
Co-chair, Open Polar Data Working Group & Tideglass Research Cooperative & 2024 -- Present & 2024-02-01 \\
Faculty representative --- \href{/blog/2026/reading-a-core/}{field note} & Pacific research partnerships committee & 2022 -- 2024 \\%
}

\newcommand{\cvLeadershipHeader}{%
\textbf{Title} & \textbf{Org} & \textbf{Dates} & \textbf{Announced} \\%
}

\newcommand{\cvLeadershipInline}{%
Co-chair, Open Polar Data Working Group, Faculty representative --- \href{/blog/2026/reading-a-core/}{field note}%
}

\newcommand{\cvLeadershipCount}{2}

\newcommand{\cvLeadershipPrinted}{0}

\newcommand{\cvLanguagesNote}{%
%
}

\newcommand{\cvLanguages}{%
\ifnum1>\cvmax\else
\cventry
  {Tamil}
  {Native}
  {}
  {}
\fi

\ifnum2>\cvmax\else
\cventry
  {English}
  {Fluent}
  {}
  {}
\fi

\ifnum3>\cvmax\else
\cventry
  {te reo Māori}
  {Intermediate learner}
  {}
  {}
\fi%
}

\newcommand{\cvLanguagesRows}{%
Tamil & Native \\
English & Fluent \\
te reo Māori & Intermediate learner \\%
}

\newcommand{\cvLanguagesHeader}{%
\textbf{Title} & \textbf{Detail} \\%
}

\newcommand{\cvLanguagesInline}{%
Tamil (Native), English (Fluent), te reo Māori (Intermediate learner)%
}

\newcommand{\cvLanguagesCount}{3}

\newcommand{\cvPublicationsCount}{10}

\defbibfilter{Publications1}{type=article and not keyword=underreview}
\defbibfilter{Publications2}{type=inproceedings and not ( type=article and not keyword=underreview )}
\defbibfilter{Publications3}{( type=incollection or type=inbook or type=book ) and not ( type=article and not keyword=underreview ) and not ( type=inproceedings )}
\defbibfilter{Publications4}{type=dataset and not ( type=article and not keyword=underreview ) and not ( type=inproceedings ) and not ( ( type=incollection or type=inbook or type=book ) )}
\defbibfilter{Publications6}{type=article and keyword=underreview and not ( type=article and not keyword=underreview ) and not ( type=inproceedings ) and not ( ( type=incollection or type=inbook or type=book ) ) and not ( type=dataset ) and not ( type=misc )}

\newcommand{\cvPublicationsKey}{%
J=Journal, C=Conference, B=Books, D=Data, R=Under review%
}

\newcommand{\cvPublicationsSections}{%
\begin{refcontext}[labelprefix=J]
\printbibliography[heading=bibsubheading, title={Journal articles (peer-reviewed)}, filter=Publications1, resetnumbers=true]
\end{refcontext}

\begin{refcontext}[labelprefix=C]
\printbibliography[heading=bibsubheading, title={Conference papers (peer-reviewed)}, filter=Publications2, resetnumbers=true]
\end{refcontext}

\begin{refcontext}[labelprefix=B]
\printbibliography[heading=bibsubheading, title={Books and chapters}, filter=Publications3, resetnumbers=true]
\end{refcontext}

\begin{refcontext}[labelprefix=D]
\printbibliography[heading=bibsubheading, title={Data archives}, filter=Publications4, resetnumbers=true]
\end{refcontext}

\begin{refcontext}[labelprefix=R]
\printbibliography[heading=bibsubheading, title={Manuscripts under review}, filter=Publications6, resetnumbers=true]
\end{refcontext}%
}

\newcommand{\cvTalksCount}{4}

\defbibfilter{Talks1}{type=unpublished and keyword=invited}

\newcommand{\cvTalksKey}{%
I=Invited%
}

\newcommand{\cvTalksSections}{%
\begin{refcontext}[labelprefix=I]
\printbibliography[heading=bibsubheading, title={Invited and keynote talks}, filter=Talks1, resetnumbers=true]
\end{refcontext}%
}

\newcommand{\cvAutoSections}{%
\cvautopart{Teaching Statement}{TeachingStatement}%
\cvautopart{Strands}{Strands}%
\cvautopart{Appointments}{Appointments}%
\cvautopart{Education}{Education}%
\cvautopart{Teaching}{Teaching}%
\cvautopart{Supervision}{Supervision}%
\cvautopart{Service}{Service}%
\cvautopart{Projects}{Projects}%
\cvautopart{Fieldwork}{Fieldwork}%
\cvautopart{Public engagement}{Outreach}%
\cvautopart{Awards}{Awards}%
\cvautopart{Leadership}{Leadership}%
\cvautopart{Languages}{Languages}%
}


